\documentclass[11pt]{article}

    \usepackage[breakable]{tcolorbox}
    \usepackage{parskip} % Stop auto-indenting (to mimic markdown behaviour)
    

    % Basic figure setup, for now with no caption control since it's done
    % automatically by Pandoc (which extracts ![](path) syntax from Markdown).
    \usepackage{graphicx}
    % Keep aspect ratio if custom image width or height is specified
    \setkeys{Gin}{keepaspectratio}
    % Maintain compatibility with old templates. Remove in nbconvert 6.0
    \let\Oldincludegraphics\includegraphics
    % Ensure that by default, figures have no caption (until we provide a
    % proper Figure object with a Caption API and a way to capture that
    % in the conversion process - todo).
    \usepackage{caption}
    \DeclareCaptionFormat{nocaption}{}
    \captionsetup{format=nocaption,aboveskip=0pt,belowskip=0pt}

    \usepackage{float}
    \floatplacement{figure}{H} % forces figures to be placed at the correct location
    \usepackage{xcolor} % Allow colors to be defined
    \usepackage{enumerate} % Needed for markdown enumerations to work
    \usepackage{geometry} % Used to adjust the document margins
    \usepackage{amsmath} % Equations
    \usepackage{amssymb} % Equations
    \usepackage{textcomp} % defines textquotesingle
    % Hack from http://tex.stackexchange.com/a/47451/13684:
    \AtBeginDocument{%
        \def\PYZsq{\textquotesingle}% Upright quotes in Pygmentized code
    }
    \usepackage{upquote} % Upright quotes for verbatim code
    \usepackage{eurosym} % defines \euro

    \usepackage{iftex}
    \ifPDFTeX
        \usepackage[T1]{fontenc}
        \IfFileExists{alphabeta.sty}{
              \usepackage{alphabeta}
          }{
              \usepackage[mathletters]{ucs}
              \usepackage[utf8x]{inputenc}
          }
    \else
        \usepackage{fontspec}
        \usepackage{unicode-math}
    \fi

    \usepackage{fancyvrb} % verbatim replacement that allows latex
    \usepackage{grffile} % extends the file name processing of package graphics
                         % to support a larger range
    \makeatletter % fix for old versions of grffile with XeLaTeX
    \@ifpackagelater{grffile}{2019/11/01}
    {
      % Do nothing on new versions
    }
    {
      \def\Gread@@xetex#1{%
        \IfFileExists{"\Gin@base".bb}%
        {\Gread@eps{\Gin@base.bb}}%
        {\Gread@@xetex@aux#1}%
      }
    }
    \makeatother
    \usepackage[Export]{adjustbox} % Used to constrain images to a maximum size
    \adjustboxset{max size={0.9\linewidth}{0.9\paperheight}}

    % The hyperref package gives us a pdf with properly built
    % internal navigation ('pdf bookmarks' for the table of contents,
    % internal cross-reference links, web links for URLs, etc.)
    \usepackage{hyperref}
    % The default LaTeX title has an obnoxious amount of whitespace. By default,
    % titling removes some of it. It also provides customization options.
    \usepackage{titling}
    \usepackage{longtable} % longtable support required by pandoc >1.10
    \usepackage{booktabs}  % table support for pandoc > 1.12.2
    \usepackage{array}     % table support for pandoc >= 2.11.3
    \usepackage{calc}      % table minipage width calculation for pandoc >= 2.11.1
    \usepackage[inline]{enumitem} % IRkernel/repr support (it uses the enumerate* environment)
    \usepackage[normalem]{ulem} % ulem is needed to support strikethroughs (\sout)
                                % normalem makes italics be italics, not underlines
    \usepackage{soul}      % strikethrough (\st) support for pandoc >= 3.0.0
    \usepackage{mathrsfs}
    

    
    % Colors for the hyperref package
    \definecolor{urlcolor}{rgb}{0,.145,.698}
    \definecolor{linkcolor}{rgb}{.71,0.21,0.01}
    \definecolor{citecolor}{rgb}{.12,.54,.11}

    % ANSI colors
    \definecolor{ansi-black}{HTML}{3E424D}
    \definecolor{ansi-black-intense}{HTML}{282C36}
    \definecolor{ansi-red}{HTML}{E75C58}
    \definecolor{ansi-red-intense}{HTML}{B22B31}
    \definecolor{ansi-green}{HTML}{00A250}
    \definecolor{ansi-green-intense}{HTML}{007427}
    \definecolor{ansi-yellow}{HTML}{DDB62B}
    \definecolor{ansi-yellow-intense}{HTML}{B27D12}
    \definecolor{ansi-blue}{HTML}{208FFB}
    \definecolor{ansi-blue-intense}{HTML}{0065CA}
    \definecolor{ansi-magenta}{HTML}{D160C4}
    \definecolor{ansi-magenta-intense}{HTML}{A03196}
    \definecolor{ansi-cyan}{HTML}{60C6C8}
    \definecolor{ansi-cyan-intense}{HTML}{258F8F}
    \definecolor{ansi-white}{HTML}{C5C1B4}
    \definecolor{ansi-white-intense}{HTML}{A1A6B2}
    \definecolor{ansi-default-inverse-fg}{HTML}{FFFFFF}
    \definecolor{ansi-default-inverse-bg}{HTML}{000000}

    % common color for the border for error outputs.
    \definecolor{outerrorbackground}{HTML}{FFDFDF}

    % commands and environments needed by pandoc snippets
    % extracted from the output of `pandoc -s`
    \providecommand{\tightlist}{%
      \setlength{\itemsep}{0pt}\setlength{\parskip}{0pt}}
    \DefineVerbatimEnvironment{Highlighting}{Verbatim}{commandchars=\\\{\}}
    % Add ',fontsize=\small' for more characters per line
    \newenvironment{Shaded}{}{}
    \newcommand{\KeywordTok}[1]{\textcolor[rgb]{0.00,0.44,0.13}{\textbf{{#1}}}}
    \newcommand{\DataTypeTok}[1]{\textcolor[rgb]{0.56,0.13,0.00}{{#1}}}
    \newcommand{\DecValTok}[1]{\textcolor[rgb]{0.25,0.63,0.44}{{#1}}}
    \newcommand{\BaseNTok}[1]{\textcolor[rgb]{0.25,0.63,0.44}{{#1}}}
    \newcommand{\FloatTok}[1]{\textcolor[rgb]{0.25,0.63,0.44}{{#1}}}
    \newcommand{\CharTok}[1]{\textcolor[rgb]{0.25,0.44,0.63}{{#1}}}
    \newcommand{\StringTok}[1]{\textcolor[rgb]{0.25,0.44,0.63}{{#1}}}
    \newcommand{\CommentTok}[1]{\textcolor[rgb]{0.38,0.63,0.69}{\textit{{#1}}}}
    \newcommand{\OtherTok}[1]{\textcolor[rgb]{0.00,0.44,0.13}{{#1}}}
    \newcommand{\AlertTok}[1]{\textcolor[rgb]{1.00,0.00,0.00}{\textbf{{#1}}}}
    \newcommand{\FunctionTok}[1]{\textcolor[rgb]{0.02,0.16,0.49}{{#1}}}
    \newcommand{\RegionMarkerTok}[1]{{#1}}
    \newcommand{\ErrorTok}[1]{\textcolor[rgb]{1.00,0.00,0.00}{\textbf{{#1}}}}
    \newcommand{\NormalTok}[1]{{#1}}

    % Additional commands for more recent versions of Pandoc
    \newcommand{\ConstantTok}[1]{\textcolor[rgb]{0.53,0.00,0.00}{{#1}}}
    \newcommand{\SpecialCharTok}[1]{\textcolor[rgb]{0.25,0.44,0.63}{{#1}}}
    \newcommand{\VerbatimStringTok}[1]{\textcolor[rgb]{0.25,0.44,0.63}{{#1}}}
    \newcommand{\SpecialStringTok}[1]{\textcolor[rgb]{0.73,0.40,0.53}{{#1}}}
    \newcommand{\ImportTok}[1]{{#1}}
    \newcommand{\DocumentationTok}[1]{\textcolor[rgb]{0.73,0.13,0.13}{\textit{{#1}}}}
    \newcommand{\AnnotationTok}[1]{\textcolor[rgb]{0.38,0.63,0.69}{\textbf{\textit{{#1}}}}}
    \newcommand{\CommentVarTok}[1]{\textcolor[rgb]{0.38,0.63,0.69}{\textbf{\textit{{#1}}}}}
    \newcommand{\VariableTok}[1]{\textcolor[rgb]{0.10,0.09,0.49}{{#1}}}
    \newcommand{\ControlFlowTok}[1]{\textcolor[rgb]{0.00,0.44,0.13}{\textbf{{#1}}}}
    \newcommand{\OperatorTok}[1]{\textcolor[rgb]{0.40,0.40,0.40}{{#1}}}
    \newcommand{\BuiltInTok}[1]{{#1}}
    \newcommand{\ExtensionTok}[1]{{#1}}
    \newcommand{\PreprocessorTok}[1]{\textcolor[rgb]{0.74,0.48,0.00}{{#1}}}
    \newcommand{\AttributeTok}[1]{\textcolor[rgb]{0.49,0.56,0.16}{{#1}}}
    \newcommand{\InformationTok}[1]{\textcolor[rgb]{0.38,0.63,0.69}{\textbf{\textit{{#1}}}}}
    \newcommand{\WarningTok}[1]{\textcolor[rgb]{0.38,0.63,0.69}{\textbf{\textit{{#1}}}}}
    \makeatletter
    \newsavebox\pandoc@box
    \newcommand*\pandocbounded[1]{%
      \sbox\pandoc@box{#1}%
      % scaling factors for width and height
      \Gscale@div\@tempa\textheight{\dimexpr\ht\pandoc@box+\dp\pandoc@box\relax}%
      \Gscale@div\@tempb\linewidth{\wd\pandoc@box}%
      % select the smaller of both
      \ifdim\@tempb\p@<\@tempa\p@
        \let\@tempa\@tempb
      \fi
      % scaling accordingly (\@tempa < 1)
      \ifdim\@tempa\p@<\p@
        \scalebox{\@tempa}{\usebox\pandoc@box}%
      % scaling not needed, use as it is
      \else
        \usebox{\pandoc@box}%
      \fi
    }
    \makeatother

    % Define a nice break command that doesn't care if a line doesn't already
    % exist.
    \def\br{\hspace*{\fill} \\* }
    % Math Jax compatibility definitions
    \def\gt{>}
    \def\lt{<}
    \let\Oldtex\TeX
    \let\Oldlatex\LaTeX
    \renewcommand{\TeX}{\textrm{\Oldtex}}
    \renewcommand{\LaTeX}{\textrm{\Oldlatex}}
    % Document parameters
    % Document title
    \title{u\_ism\_and\_star\_birth}
    
    
    
    
    
    
    
% Pygments definitions
\makeatletter
\def\PY@reset{\let\PY@it=\relax \let\PY@bf=\relax%
    \let\PY@ul=\relax \let\PY@tc=\relax%
    \let\PY@bc=\relax \let\PY@ff=\relax}
\def\PY@tok#1{\csname PY@tok@#1\endcsname}
\def\PY@toks#1+{\ifx\relax#1\empty\else%
    \PY@tok{#1}\expandafter\PY@toks\fi}
\def\PY@do#1{\PY@bc{\PY@tc{\PY@ul{%
    \PY@it{\PY@bf{\PY@ff{#1}}}}}}}
\def\PY#1#2{\PY@reset\PY@toks#1+\relax+\PY@do{#2}}

\@namedef{PY@tok@w}{\def\PY@tc##1{\textcolor[rgb]{0.73,0.73,0.73}{##1}}}
\@namedef{PY@tok@c}{\let\PY@it=\textit\def\PY@tc##1{\textcolor[rgb]{0.24,0.48,0.48}{##1}}}
\@namedef{PY@tok@cp}{\def\PY@tc##1{\textcolor[rgb]{0.61,0.40,0.00}{##1}}}
\@namedef{PY@tok@k}{\let\PY@bf=\textbf\def\PY@tc##1{\textcolor[rgb]{0.00,0.50,0.00}{##1}}}
\@namedef{PY@tok@kp}{\def\PY@tc##1{\textcolor[rgb]{0.00,0.50,0.00}{##1}}}
\@namedef{PY@tok@kt}{\def\PY@tc##1{\textcolor[rgb]{0.69,0.00,0.25}{##1}}}
\@namedef{PY@tok@o}{\def\PY@tc##1{\textcolor[rgb]{0.40,0.40,0.40}{##1}}}
\@namedef{PY@tok@ow}{\let\PY@bf=\textbf\def\PY@tc##1{\textcolor[rgb]{0.67,0.13,1.00}{##1}}}
\@namedef{PY@tok@nb}{\def\PY@tc##1{\textcolor[rgb]{0.00,0.50,0.00}{##1}}}
\@namedef{PY@tok@nf}{\def\PY@tc##1{\textcolor[rgb]{0.00,0.00,1.00}{##1}}}
\@namedef{PY@tok@nc}{\let\PY@bf=\textbf\def\PY@tc##1{\textcolor[rgb]{0.00,0.00,1.00}{##1}}}
\@namedef{PY@tok@nn}{\let\PY@bf=\textbf\def\PY@tc##1{\textcolor[rgb]{0.00,0.00,1.00}{##1}}}
\@namedef{PY@tok@ne}{\let\PY@bf=\textbf\def\PY@tc##1{\textcolor[rgb]{0.80,0.25,0.22}{##1}}}
\@namedef{PY@tok@nv}{\def\PY@tc##1{\textcolor[rgb]{0.10,0.09,0.49}{##1}}}
\@namedef{PY@tok@no}{\def\PY@tc##1{\textcolor[rgb]{0.53,0.00,0.00}{##1}}}
\@namedef{PY@tok@nl}{\def\PY@tc##1{\textcolor[rgb]{0.46,0.46,0.00}{##1}}}
\@namedef{PY@tok@ni}{\let\PY@bf=\textbf\def\PY@tc##1{\textcolor[rgb]{0.44,0.44,0.44}{##1}}}
\@namedef{PY@tok@na}{\def\PY@tc##1{\textcolor[rgb]{0.41,0.47,0.13}{##1}}}
\@namedef{PY@tok@nt}{\let\PY@bf=\textbf\def\PY@tc##1{\textcolor[rgb]{0.00,0.50,0.00}{##1}}}
\@namedef{PY@tok@nd}{\def\PY@tc##1{\textcolor[rgb]{0.67,0.13,1.00}{##1}}}
\@namedef{PY@tok@s}{\def\PY@tc##1{\textcolor[rgb]{0.73,0.13,0.13}{##1}}}
\@namedef{PY@tok@sd}{\let\PY@it=\textit\def\PY@tc##1{\textcolor[rgb]{0.73,0.13,0.13}{##1}}}
\@namedef{PY@tok@si}{\let\PY@bf=\textbf\def\PY@tc##1{\textcolor[rgb]{0.64,0.35,0.47}{##1}}}
\@namedef{PY@tok@se}{\let\PY@bf=\textbf\def\PY@tc##1{\textcolor[rgb]{0.67,0.36,0.12}{##1}}}
\@namedef{PY@tok@sr}{\def\PY@tc##1{\textcolor[rgb]{0.64,0.35,0.47}{##1}}}
\@namedef{PY@tok@ss}{\def\PY@tc##1{\textcolor[rgb]{0.10,0.09,0.49}{##1}}}
\@namedef{PY@tok@sx}{\def\PY@tc##1{\textcolor[rgb]{0.00,0.50,0.00}{##1}}}
\@namedef{PY@tok@m}{\def\PY@tc##1{\textcolor[rgb]{0.40,0.40,0.40}{##1}}}
\@namedef{PY@tok@gh}{\let\PY@bf=\textbf\def\PY@tc##1{\textcolor[rgb]{0.00,0.00,0.50}{##1}}}
\@namedef{PY@tok@gu}{\let\PY@bf=\textbf\def\PY@tc##1{\textcolor[rgb]{0.50,0.00,0.50}{##1}}}
\@namedef{PY@tok@gd}{\def\PY@tc##1{\textcolor[rgb]{0.63,0.00,0.00}{##1}}}
\@namedef{PY@tok@gi}{\def\PY@tc##1{\textcolor[rgb]{0.00,0.52,0.00}{##1}}}
\@namedef{PY@tok@gr}{\def\PY@tc##1{\textcolor[rgb]{0.89,0.00,0.00}{##1}}}
\@namedef{PY@tok@ge}{\let\PY@it=\textit}
\@namedef{PY@tok@gs}{\let\PY@bf=\textbf}
\@namedef{PY@tok@ges}{\let\PY@bf=\textbf\let\PY@it=\textit}
\@namedef{PY@tok@gp}{\let\PY@bf=\textbf\def\PY@tc##1{\textcolor[rgb]{0.00,0.00,0.50}{##1}}}
\@namedef{PY@tok@go}{\def\PY@tc##1{\textcolor[rgb]{0.44,0.44,0.44}{##1}}}
\@namedef{PY@tok@gt}{\def\PY@tc##1{\textcolor[rgb]{0.00,0.27,0.87}{##1}}}
\@namedef{PY@tok@err}{\def\PY@bc##1{{\setlength{\fboxsep}{\string -\fboxrule}\fcolorbox[rgb]{1.00,0.00,0.00}{1,1,1}{\strut ##1}}}}
\@namedef{PY@tok@kc}{\let\PY@bf=\textbf\def\PY@tc##1{\textcolor[rgb]{0.00,0.50,0.00}{##1}}}
\@namedef{PY@tok@kd}{\let\PY@bf=\textbf\def\PY@tc##1{\textcolor[rgb]{0.00,0.50,0.00}{##1}}}
\@namedef{PY@tok@kn}{\let\PY@bf=\textbf\def\PY@tc##1{\textcolor[rgb]{0.00,0.50,0.00}{##1}}}
\@namedef{PY@tok@kr}{\let\PY@bf=\textbf\def\PY@tc##1{\textcolor[rgb]{0.00,0.50,0.00}{##1}}}
\@namedef{PY@tok@bp}{\def\PY@tc##1{\textcolor[rgb]{0.00,0.50,0.00}{##1}}}
\@namedef{PY@tok@fm}{\def\PY@tc##1{\textcolor[rgb]{0.00,0.00,1.00}{##1}}}
\@namedef{PY@tok@vc}{\def\PY@tc##1{\textcolor[rgb]{0.10,0.09,0.49}{##1}}}
\@namedef{PY@tok@vg}{\def\PY@tc##1{\textcolor[rgb]{0.10,0.09,0.49}{##1}}}
\@namedef{PY@tok@vi}{\def\PY@tc##1{\textcolor[rgb]{0.10,0.09,0.49}{##1}}}
\@namedef{PY@tok@vm}{\def\PY@tc##1{\textcolor[rgb]{0.10,0.09,0.49}{##1}}}
\@namedef{PY@tok@sa}{\def\PY@tc##1{\textcolor[rgb]{0.73,0.13,0.13}{##1}}}
\@namedef{PY@tok@sb}{\def\PY@tc##1{\textcolor[rgb]{0.73,0.13,0.13}{##1}}}
\@namedef{PY@tok@sc}{\def\PY@tc##1{\textcolor[rgb]{0.73,0.13,0.13}{##1}}}
\@namedef{PY@tok@dl}{\def\PY@tc##1{\textcolor[rgb]{0.73,0.13,0.13}{##1}}}
\@namedef{PY@tok@s2}{\def\PY@tc##1{\textcolor[rgb]{0.73,0.13,0.13}{##1}}}
\@namedef{PY@tok@sh}{\def\PY@tc##1{\textcolor[rgb]{0.73,0.13,0.13}{##1}}}
\@namedef{PY@tok@s1}{\def\PY@tc##1{\textcolor[rgb]{0.73,0.13,0.13}{##1}}}
\@namedef{PY@tok@mb}{\def\PY@tc##1{\textcolor[rgb]{0.40,0.40,0.40}{##1}}}
\@namedef{PY@tok@mf}{\def\PY@tc##1{\textcolor[rgb]{0.40,0.40,0.40}{##1}}}
\@namedef{PY@tok@mh}{\def\PY@tc##1{\textcolor[rgb]{0.40,0.40,0.40}{##1}}}
\@namedef{PY@tok@mi}{\def\PY@tc##1{\textcolor[rgb]{0.40,0.40,0.40}{##1}}}
\@namedef{PY@tok@il}{\def\PY@tc##1{\textcolor[rgb]{0.40,0.40,0.40}{##1}}}
\@namedef{PY@tok@mo}{\def\PY@tc##1{\textcolor[rgb]{0.40,0.40,0.40}{##1}}}
\@namedef{PY@tok@ch}{\let\PY@it=\textit\def\PY@tc##1{\textcolor[rgb]{0.24,0.48,0.48}{##1}}}
\@namedef{PY@tok@cm}{\let\PY@it=\textit\def\PY@tc##1{\textcolor[rgb]{0.24,0.48,0.48}{##1}}}
\@namedef{PY@tok@cpf}{\let\PY@it=\textit\def\PY@tc##1{\textcolor[rgb]{0.24,0.48,0.48}{##1}}}
\@namedef{PY@tok@c1}{\let\PY@it=\textit\def\PY@tc##1{\textcolor[rgb]{0.24,0.48,0.48}{##1}}}
\@namedef{PY@tok@cs}{\let\PY@it=\textit\def\PY@tc##1{\textcolor[rgb]{0.24,0.48,0.48}{##1}}}

\def\PYZbs{\char`\\}
\def\PYZus{\char`\_}
\def\PYZob{\char`\{}
\def\PYZcb{\char`\}}
\def\PYZca{\char`\^}
\def\PYZam{\char`\&}
\def\PYZlt{\char`\<}
\def\PYZgt{\char`\>}
\def\PYZsh{\char`\#}
\def\PYZpc{\char`\%}
\def\PYZdl{\char`\$}
\def\PYZhy{\char`\-}
\def\PYZsq{\char`\'}
\def\PYZdq{\char`\"}
\def\PYZti{\char`\~}
% for compatibility with earlier versions
\def\PYZat{@}
\def\PYZlb{[}
\def\PYZrb{]}
\makeatother


    % For linebreaks inside Verbatim environment from package fancyvrb.
    \makeatletter
        \newbox\Wrappedcontinuationbox
        \newbox\Wrappedvisiblespacebox
        \newcommand*\Wrappedvisiblespace {\textcolor{red}{\textvisiblespace}}
        \newcommand*\Wrappedcontinuationsymbol {\textcolor{red}{\llap{\tiny$\m@th\hookrightarrow$}}}
        \newcommand*\Wrappedcontinuationindent {3ex }
        \newcommand*\Wrappedafterbreak {\kern\Wrappedcontinuationindent\copy\Wrappedcontinuationbox}
        % Take advantage of the already applied Pygments mark-up to insert
        % potential linebreaks for TeX processing.
        %        {, <, #, %, $, ' and ": go to next line.
        %        _, }, ^, &, >, - and ~: stay at end of broken line.
        % Use of \textquotesingle for straight quote.
        \newcommand*\Wrappedbreaksatspecials {%
            \def\PYGZus{\discretionary{\char`\_}{\Wrappedafterbreak}{\char`\_}}%
            \def\PYGZob{\discretionary{}{\Wrappedafterbreak\char`\{}{\char`\{}}%
            \def\PYGZcb{\discretionary{\char`\}}{\Wrappedafterbreak}{\char`\}}}%
            \def\PYGZca{\discretionary{\char`\^}{\Wrappedafterbreak}{\char`\^}}%
            \def\PYGZam{\discretionary{\char`\&}{\Wrappedafterbreak}{\char`\&}}%
            \def\PYGZlt{\discretionary{}{\Wrappedafterbreak\char`\<}{\char`\<}}%
            \def\PYGZgt{\discretionary{\char`\>}{\Wrappedafterbreak}{\char`\>}}%
            \def\PYGZsh{\discretionary{}{\Wrappedafterbreak\char`\#}{\char`\#}}%
            \def\PYGZpc{\discretionary{}{\Wrappedafterbreak\char`\%}{\char`\%}}%
            \def\PYGZdl{\discretionary{}{\Wrappedafterbreak\char`\$}{\char`\$}}%
            \def\PYGZhy{\discretionary{\char`\-}{\Wrappedafterbreak}{\char`\-}}%
            \def\PYGZsq{\discretionary{}{\Wrappedafterbreak\textquotesingle}{\textquotesingle}}%
            \def\PYGZdq{\discretionary{}{\Wrappedafterbreak\char`\"}{\char`\"}}%
            \def\PYGZti{\discretionary{\char`\~}{\Wrappedafterbreak}{\char`\~}}%
        }
        % Some characters . , ; ? ! / are not pygmentized.
        % This macro makes them "active" and they will insert potential linebreaks
        \newcommand*\Wrappedbreaksatpunct {%
            \lccode`\~`\.\lowercase{\def~}{\discretionary{\hbox{\char`\.}}{\Wrappedafterbreak}{\hbox{\char`\.}}}%
            \lccode`\~`\,\lowercase{\def~}{\discretionary{\hbox{\char`\,}}{\Wrappedafterbreak}{\hbox{\char`\,}}}%
            \lccode`\~`\;\lowercase{\def~}{\discretionary{\hbox{\char`\;}}{\Wrappedafterbreak}{\hbox{\char`\;}}}%
            \lccode`\~`\:\lowercase{\def~}{\discretionary{\hbox{\char`\:}}{\Wrappedafterbreak}{\hbox{\char`\:}}}%
            \lccode`\~`\?\lowercase{\def~}{\discretionary{\hbox{\char`\?}}{\Wrappedafterbreak}{\hbox{\char`\?}}}%
            \lccode`\~`\!\lowercase{\def~}{\discretionary{\hbox{\char`\!}}{\Wrappedafterbreak}{\hbox{\char`\!}}}%
            \lccode`\~`\/\lowercase{\def~}{\discretionary{\hbox{\char`\/}}{\Wrappedafterbreak}{\hbox{\char`\/}}}%
            \catcode`\.\active
            \catcode`\,\active
            \catcode`\;\active
            \catcode`\:\active
            \catcode`\?\active
            \catcode`\!\active
            \catcode`\/\active
            \lccode`\~`\~
        }
    \makeatother

    \let\OriginalVerbatim=\Verbatim
    \makeatletter
    \renewcommand{\Verbatim}[1][1]{%
        %\parskip\z@skip
        \sbox\Wrappedcontinuationbox {\Wrappedcontinuationsymbol}%
        \sbox\Wrappedvisiblespacebox {\FV@SetupFont\Wrappedvisiblespace}%
        \def\FancyVerbFormatLine ##1{\hsize\linewidth
            \vtop{\raggedright\hyphenpenalty\z@\exhyphenpenalty\z@
                \doublehyphendemerits\z@\finalhyphendemerits\z@
                \strut ##1\strut}%
        }%
        % If the linebreak is at a space, the latter will be displayed as visible
        % space at end of first line, and a continuation symbol starts next line.
        % Stretch/shrink are however usually zero for typewriter font.
        \def\FV@Space {%
            \nobreak\hskip\z@ plus\fontdimen3\font minus\fontdimen4\font
            \discretionary{\copy\Wrappedvisiblespacebox}{\Wrappedafterbreak}
            {\kern\fontdimen2\font}%
        }%

        % Allow breaks at special characters using \PYG... macros.
        \Wrappedbreaksatspecials
        % Breaks at punctuation characters . , ; ? ! and / need catcode=\active
        \OriginalVerbatim[#1,codes*=\Wrappedbreaksatpunct]%
    }
    \makeatother

    % Exact colors from NB
    \definecolor{incolor}{HTML}{303F9F}
    \definecolor{outcolor}{HTML}{D84315}
    \definecolor{cellborder}{HTML}{CFCFCF}
    \definecolor{cellbackground}{HTML}{F7F7F7}

    % prompt
    \makeatletter
    \newcommand{\boxspacing}{\kern\kvtcb@left@rule\kern\kvtcb@boxsep}
    \makeatother
    \newcommand{\prompt}[4]{
        {\ttfamily\llap{{\color{#2}[#3]:\hspace{3pt}#4}}\vspace{-\baselineskip}}
    }
    

    
    % Prevent overflowing lines due to hard-to-break entities
    \sloppy
    % Setup hyperref package
    \hypersetup{
      breaklinks=true,  % so long urls are correctly broken across lines
      colorlinks=true,
      urlcolor=urlcolor,
      linkcolor=linkcolor,
      citecolor=citecolor,
      }
    % Slightly bigger margins than the latex defaults
    
    \geometry{verbose,tmargin=1in,bmargin=1in,lmargin=1in,rmargin=1in}
    
    

\begin{document}
    
    \maketitle
    
    

    
    \section{\texorpdfstring{\href{https://en.wikipedia.org/wiki/Interstellar_medium}{ISM}
and \href{https://en.wikipedia.org/wiki/Star_formation}{Star
Formation}}{ISM and Star Formation}}\label{ism-and-star-formation}

In this packet, we turn to the study of stellar formation. The details
of stellar formation and evolution depend sensitively on masses and
compositions. To get oriented, we start with the formation of stars that
are very similar to our Sun, after which, we can turn to stars that are
much bigger, smaller, older, past their primes, \emph{et cetera}.

 
Throughout this packet, we will assume that the stars we are talking
about have masses and metalicities similar to the Sun.
 

In AA:SS, we studied the formation of the Sun from a fragment of a Giant
Molecular Cloud (GMC). Rather than redo the whole thing for very similar
stars, we will review some aspects of that process only very briefly and
focus on other aspects that we did not emphasize in
\texttt{ss\_solar\_system\_1,2,3}. For students who did not take AA:SS
and would like to know more, I recommend you look over those packets
and/or take AA:SS next year. For students who did take AA:SS, I
recommend you review your notes from that class. Of course, I am always
happy to answer any questions you may have about * homologous cloud
collapse and fragmentation * core heating throughout stages * accretion
mechanisms * anything else

\subsection{\texorpdfstring{\href{https://en.wikipedia.org/wiki/Interstellar_medium}{Interstellar
Matter}}{Interstellar Matter}}\label{interstellar-matter}

In space we notice the stars, and to a lesser extent the planets, but
there is \emph{much} more material there that we don't see. We call this
\emph{interstellar matter} or ISM. In fact, we do ``see'' ISM when we
know how to look. ISM is roughly divided into three categories:

\begin{itemize}
\tightlist
\item
  atomic (size \textasciitilde{} 1 AA \(=10^{-10}\) m = 0.1 nm)\\
\item
  molecular (size \textasciitilde{} 10 AA \(=10^{-9}\) m = 1 nm)\\
\item
  dust (size \textasciitilde{} 1000 AA \(=10^{-7}\) m = 0.1 μm)
\end{itemize}

The dust does not emit light, but it reflects/absorbs light making for
dark spots in the sky where otherwise there would be light. The
interactions of particles with light are collectively called
\textbf{scattering} and the processes through which this happens are
called \textbf{scattering} processes. The diminution of light by dust is
called
\href{https://en.wikipedia.org/wiki/Extinction_(astronomy)}{\textbf{extinction}}
and \textbf{reddening}. In the first effect, light is re-reflected away
from the line of sight and lost into space. Particles with sizes \(d\)
much larger than the wavelength of the light \(\lambda\) simply block
the light and cast a shadow with an area denoted by \(\sigma\) called
the \textbf{cross section} of the particle.

\subsubsection{\texorpdfstring{\href{https://en.wikipedia.org/wiki/Rayleigh_scattering}{Rayleigh
scattering}}{Rayleigh scattering}}\label{rayleigh-scattering}

is the is the scattering of light by particles with size
\(d\ll \lambda\) is much smaller than the wavelength of the light. The
details are complicated but the formula for the intensity of the
scattered light is \[
    I = c I_0 \left(\frac{1+\cos^2\theta}2\right) \frac {d^6}{r^2}\frac1{\lambda^4} \quad \quad \quad(\lambda \gg d)    \tag{1}
\] This is called the Rayleigh formula for scattering. Here \(c\) is
some dimensionless constant, \(I_0\) is the intensity of the light
before it is scattered, \(r\) is the distance to the particle doing the
scattering (note the characteristic \(r^{-2}\) fall-off!), and
\(\theta\) denotes the direction in which the light is scattered. This
formula says that light is scattered in all directions and says what
fraction goes in each direction.

 
\subsubsection{Problem 1}\label{problem-1}

\begin{enumerate}
\def\labelenumi{\alph{enumi}.}
\item
  Explain why the radial dependence for the intensity falls off like
  \(1/r^2\). (Hint: this is the fall-off of intensity with distance
  we've discussed before.)
\item
  Plot the \(\frac12\left(1+\cos^2\theta\right)\) factor in Desmos to
  see what the radiation field looks like around the scattering
  particle.
\item
  Note the strong dependence on the wavelength and use it to explain why
  the sun is white/yellow, the sky is blue, and sunsets are red.
  Calculate a range for how much more strongly blue light is scattered
  compared to red.
\item
  For visible light, to what type of interstellar matter does the
  Rayleigh solution apply?
\end{enumerate}
 

     
\subsubsection{Answer 1}\label{answer-1}

\begin{enumerate}
\def\labelenumi{\alph{enumi}.}
\item
  This is the usual \(1/r^2\) fall-off of any massless force carrier. We
  saw this behavior for both the gravitational force law and the Coulomb
  force.
\item
  The angular distribution of Rayleigh-scattered radiation. The curve
  has been exaggerated to \((1+a \cos^2\theta)\) with \(a=4\) to
  emphasize the double-lobe structure.
\item
  The characteristic \(\lambda^{-4}\) dependence implies that
  short-wavelength light is affected far more strongly than
  long-wavelength light. For example, the visible spectrum runs from 380
  \texttt{nm} (violet) to 750 \texttt{nm} (red), so violet is scattered
  \(\left(\frac{750}{380}\right)^4  \left(1.97\right)^4=15\) times more
  than red.
\end{enumerate}

If we looked directly at the Sun (don't do it!!) it would look white,
but if we look off to some random point in the clear sky, it would look
blue, because this is the color that has the best chance of being
scattered back to us. During sunsets, on the other hand, we generally
look somewhere near the direction of the Sun. Now we are looking through
a lot of atmosphere so the Sun is no longer white, but red, since the
blue light was mostly scattered out on it's way to our eyes.

Below, we compare the scattering of blue and red light quantitiatively.

\begin{enumerate}
\def\labelenumi{\alph{enumi}.}
\setcounter{enumi}{3}
\tightlist
\item
  For visible light,
  \(\lambda \sim 500\ nm = 5\times 10^{-7}\ m \gg d\), so \(d\) must be
  the size of atoms or molecules. Dust grains are roughly 5 times
  smaller than this, so a different approximation would be more
  appropriate for them, however, the applicable
  \href{https://en.wikipedia.org/wiki/Mie_scattering}{Mie theory} is far
  more complicated than we should get into here.
\end{enumerate}
 

    \begin{tcolorbox}[breakable, size=fbox, boxrule=1pt, pad at break*=1mm,colback=cellbackground, colframe=cellborder]
\prompt{In}{incolor}{1}{\boxspacing}
\begin{Verbatim}[commandchars=\\\{\}]
\PY{k+kn}{import}\PY{+w}{ }\PY{n+nn}{astropy}\PY{n+nn}{.}\PY{n+nn}{units}\PY{+w}{ }\PY{k}{as}\PY{+w}{ }\PY{n+nn}{u}

\PY{n+nb}{print}\PY{p}{(}\PY{l+s+sa}{f}\PY{l+s+s2}{\PYZdq{}}\PY{l+s+se}{\PYZbs{}n}\PY{l+s+s2}{|   ANSWER 1}\PY{l+s+se}{\PYZbs{}n}\PY{l+s+s2}{|}\PY{l+s+s2}{\PYZdq{}}\PY{p}{)}

\PY{n}{blue\PYZus{}low} \PY{o}{=} \PY{l+m+mi}{450}\PY{o}{*}\PY{n}{u}\PY{o}{.}\PY{n}{nm}
\PY{n}{blue\PYZus{}high} \PY{o}{=} \PY{l+m+mi}{485}\PY{o}{*}\PY{n}{u}\PY{o}{.}\PY{n}{nm}
\PY{n}{red\PYZus{}low} \PY{o}{=} \PY{l+m+mi}{635}\PY{o}{*}\PY{n}{u}\PY{o}{.}\PY{n}{nm}
\PY{n}{red\PYZus{}high} \PY{o}{=} \PY{l+m+mi}{750}\PY{o}{*}\PY{n}{u}\PY{o}{.}\PY{n}{nm}

\PY{n+nb}{print}\PY{p}{(}\PY{l+s+sa}{f}\PY{l+s+s2}{\PYZdq{}}\PY{l+s+s2}{|   1c. For blue light between }\PY{l+s+si}{\PYZob{}}\PY{n}{blue\PYZus{}low}\PY{l+s+si}{\PYZcb{}}\PY{l+s+s2}{\PYZhy{}}\PY{l+s+si}{\PYZob{}}\PY{n}{blue\PYZus{}high}\PY{l+s+si}{\PYZcb{}}\PY{l+s+s2}{ and red between }\PY{l+s+si}{\PYZob{}}\PY{n}{red\PYZus{}low}\PY{l+s+si}{\PYZcb{}}\PY{l+s+s2}{\PYZhy{}}\PY{l+s+si}{\PYZob{}}\PY{n}{red\PYZus{}high}\PY{l+s+si}{\PYZcb{}}\PY{l+s+s2}{,}\PY{l+s+se}{\PYZbs{}n}\PY{l+s+s2}{|       the scattering is anywhere between }\PY{l+s+si}{\PYZob{}}\PY{p}{(}\PY{n}{red\PYZus{}low}\PY{o}{/}\PY{n}{blue\PYZus{}high}\PY{p}{)}\PY{o}{*}\PY{o}{*}\PY{l+m+mi}{4}\PY{l+s+si}{:}\PY{l+s+s2}{.0f}\PY{l+s+si}{\PYZcb{}}\PY{l+s+s2}{ and }\PY{l+s+si}{\PYZob{}}\PY{p}{(}\PY{n}{red\PYZus{}high}\PY{o}{/}\PY{n}{blue\PYZus{}low}\PY{p}{)}\PY{o}{*}\PY{o}{*}\PY{l+m+mi}{4}\PY{l+s+si}{:}\PY{l+s+s2}{.0f}\PY{l+s+si}{\PYZcb{}}\PY{l+s+s2}{ times as strong.}\PY{l+s+se}{\PYZbs{}n}\PY{l+s+s2}{\PYZdq{}}\PY{p}{)}
\end{Verbatim}
\end{tcolorbox}

    \begin{Verbatim}[commandchars=\\\{\}]

|   ANSWER 1
|
|   1c. For blue light between 450.0 nm-485.0 nm and red between 635.0 nm-750.0
nm,
|       the scattering is anywhere between 3 and 8 times as strong.

    \end{Verbatim}

    Along these lines, shorter-wavelength light is scattered more and longer
wavelength light is scattered less. Thus, the light that has passed
through an interstellar cloud has had more blue light removed
(reradiated in all directions) and the remainder is reddish, just like
it is for sunsets on Earth.

 
\subsubsection{Problem 2}\label{problem-2}

Meet \href{https://en.wikipedia.org/wiki/Barnard_68}{Barnard 68}. It is
a cloud of dust and molecules 125 pcs away in the constellation
\href{https://en.wikipedia.org/wiki/Ophiuchus}{Ophiuchus}. Explain why
there are no stars in this hole and why those around the edges of the
hole are red.
 

     
\subsubsection{Answer 2}\label{answer-2}

A thick dust cloud is blocking (scattering) the visible light. The
tell-tale sign of this is the reddening around the edge where the dust
cloud thins and some of the longer-wavelength visible light makes it
through. The shorter-wavelength visible light is still scattered away,
making the color red.
 

    As we discussed in class, there are three main types of interstellar
cloud depending on what the hydrogen is doing. Hydrogen comes in three
forms:

\begin{itemize}
\item
  H I: Neutral hydrogen atom consisting of a proton nucleus and an
  electron bound to it in one of infinitely many orbitals indicated by
  the ``principal quantum number'' \(n=1, 2, 3, \dots\). These orbitals
  have definite energy given by the formula \[
    E_n = -13.6\ \texttt{eV} \frac 1{n^2}\quad \quad \quad n\in {\bf N} \tag{2}
  \] The \(-\) sign indicates that the electron is bound to the proton.
  (If the energy of an electron is positive, it is unbound and will
  escape the proton.) We see that as we increase \(n\to \infty\), the
  electron is getting closer and closer to freedom as \(E_n\to 0\).
  Thus, there are infinitely many bound orbitals but any energy in
  excess of 13.6 eV is enough to kick an electron in any orbital out of
  the atom. When that happens, we get
\item
  H II: Ionized hydrogen. This is a hydrogen ion that is missing its
  electron, that is, a prosaic term for the word ``proton''.
\end{itemize}

Going the other way, a hydrogen atom can bind to another hydrogen atoms
to make

\begin{itemize}
\tightlist
\item
  H2 : Molecular hydrogen. This is two hydrogen atoms bound together
  covalently by the valence electrons of each HI. The covalent binding
  energy is about 4.5 eV.
\end{itemize}

An aside on notation: Spectroscopists refer to various states of
ionization of an element, let's call it ``Em'' with a roman numeral: Em
I is the neutral version of the element Em, Em II is the singly-ionized
one (so it has charge \(+1\)), Em III is doubly-ionized, and so on.
These ionization states can be identified by the changes in their
spectral lines.

 
\textbf{WARNING:} These conventions are \emph{not} the same as those you
learned in chemistry! Make sure you know the difference, because this is
\emph{astrophysics} class.
 

 
\subsubsection{Problem 3}\label{problem-3}

\begin{enumerate}
\def\labelenumi{\alph{enumi}.}
\item
  Fe IX and Fe X are ions present in the sun's atmosphere. What are the
  charges of these ions?
\item
  What's wrong with the ions H III and He IV?
\end{enumerate}
 

     
\subsubsection{Answer 3}\label{answer-3}

\begin{enumerate}
\def\labelenumi{\alph{enumi}.}
\item
  X is the Roman symbol for 10 and IX makes it one fewer, so the charges
  are \(+8\) for FeIX and \(+9\) for FeX.
\item
  H only has one electron, so HIII would have \(-1\) electrons, which is
  impossible. Similarly, HeIV would have 3 electrons removed, but it
  only has at most 2 in the first place.
\end{enumerate}
 

    The major discoveries of the Quantum Mechanics revolution included the
realization that particles and waves are two sides of the same coin: Yes
an electron is a particle, but it also has wave-like properties like
(\href{https://en.wikiversity.org/wiki/De_Broglie_wavelength}{de
Broglie}) wavelength. Yes, light is a wave, but it comes in discrete
packets called ``photons'', so it is also a particle with a definite
energy given by the
\href{https://en.wikipedia.org/wiki/Planck_relation}{Einstein-Planck
relation} 
$$
    E=hf\quad \quad \textrm{with}\quad \quad h=6.6\times 10^{-34}\ J\cdot s  
$$
A particle of light has a very specific speed
\(c=3.0\times 10^8\ m/s\) when it travels in vacuum. This gives a
relation for the wavelength in terms of the frequency \[
    \lambda = \frac cf      \tag{4}
\]

 
\subsubsection{Problem 4}\label{problem-4}

Use these relations to calculate the wavelength of light

\begin{enumerate}
\def\labelenumi{\alph{enumi}.}
\item
  needed to dissociate the hydrogen molecule into two hydrogen atoms
\item
  needed to ionize a hydrogen atom
\item
  least energetic photon that results from a transition to the ground
  state \(n=1\)
\item
  the H\(\alpha\) line. This is the first Balmer spectral line, which
  means that is the least energetic transition to the first excited
  state \(n=2\).
\item
  Identify which part of the
  \href{https://en.wikipedia.org/wiki/Electromagnetic_spectrum}{electromagnetic
  spectrum} these photons correspond to (\emph{i.e.} visible, UV, IR,
  \emph{etc.})
\end{enumerate}

(Recall that one
\href{https://en.wikipedia.org/wiki/Electronvolt}{electronvolt} is
\(1.6\times 10^{-19}\) \texttt{J}.)
 

    \begin{tcolorbox}[breakable, size=fbox, boxrule=1pt, pad at break*=1mm,colback=cellbackground, colframe=cellborder]
\prompt{In}{incolor}{2}{\boxspacing}
\begin{Verbatim}[commandchars=\\\{\}]
\PY{k+kn}{from}\PY{+w}{ }\PY{n+nn}{astropy}\PY{n+nn}{.}\PY{n+nn}{constants}\PY{+w}{ }\PY{k+kn}{import} \PY{n}{c}\PY{p}{,} \PY{n}{h}
\PY{k+kn}{from}\PY{+w}{ }\PY{n+nn}{aa\PYZus{}tools}\PY{n+nn}{.}\PY{n+nn}{astroconst}\PY{+w}{ }\PY{k+kn}{import} \PY{n}{spectrum\PYZus{}wavelength}

\PY{n+nb}{print}\PY{p}{(}\PY{l+s+sa}{f}\PY{l+s+s2}{\PYZdq{}}\PY{l+s+se}{\PYZbs{}n}\PY{l+s+s2}{|   ANSWER 4}\PY{l+s+se}{\PYZbs{}n}\PY{l+s+s2}{|}\PY{l+s+s2}{\PYZdq{}}\PY{p}{)}

\PY{k}{def}\PY{+w}{ }\PY{n+nf}{wavelength} \PY{p}{(}\PY{n}{dE}\PY{p}{:} \PY{n}{u}\PY{o}{.}\PY{n}{Quantity}\PY{p}{)} \PY{o}{\PYZhy{}}\PY{o}{\PYZgt{}} \PY{n}{u}\PY{o}{.}\PY{n}{Quantity}\PY{p}{:}
    \PY{n}{\PYZus{}lambda} \PY{o}{=} \PY{n}{h}\PY{o}{*}\PY{n}{c}\PY{o}{/}\PY{n}{dE}
    \PY{k}{return} \PY{n}{\PYZus{}lambda}\PY{o}{.}\PY{n}{decompose}\PY{p}{(}\PY{p}{)}

\PY{n}{dE\PYZus{}covalent} \PY{o}{=} \PY{l+m+mf}{4.5} \PY{o}{*} \PY{n}{u}\PY{o}{.}\PY{n}{eV}
\PY{n}{dE\PYZus{}H} \PY{o}{=} \PY{l+m+mf}{13.6} \PY{o}{*} \PY{n}{u}\PY{o}{.}\PY{n}{eV}
\PY{n}{dE\PYZus{}12} \PY{o}{=} \PY{n}{dE\PYZus{}H}\PY{o}{*}\PY{p}{(}\PY{l+m+mi}{1}\PY{o}{\PYZhy{}}\PY{l+m+mi}{1}\PY{o}{/}\PY{l+m+mi}{2}\PY{o}{*}\PY{o}{*}\PY{l+m+mi}{2}\PY{p}{)}
\PY{n}{dE\PYZus{}23} \PY{o}{=} \PY{n}{dE\PYZus{}H}\PY{o}{*}\PY{p}{(}\PY{l+m+mi}{1}\PY{o}{/}\PY{l+m+mi}{2}\PY{o}{*}\PY{o}{*}\PY{l+m+mi}{2}\PY{o}{\PYZhy{}}\PY{l+m+mi}{1}\PY{o}{/}\PY{l+m+mi}{3}\PY{o}{*}\PY{o}{*}\PY{l+m+mi}{2}\PY{p}{)}

\PY{n+nb}{print}\PY{p}{(}\PY{l+s+sa}{f}\PY{l+s+s2}{\PYZdq{}}\PY{l+s+s2}{|   a. Recall that the molecular hydrogen covalent bond has an energy of}\PY{l+s+se}{\PYZbs{}n}\PY{l+s+s2}{|      }\PY{l+s+si}{\PYZob{}}\PY{n}{dE\PYZus{}covalent}\PY{l+s+si}{\PYZcb{}}\PY{l+s+s2}{ = }\PY{l+s+si}{\PYZob{}}\PY{n}{dE\PYZus{}covalent}\PY{o}{.}\PY{n}{to}\PY{p}{(}\PY{n}{u}\PY{o}{.}\PY{n}{J}\PY{p}{)}\PY{l+s+si}{:}\PY{l+s+s2}{.1e}\PY{l+s+si}{\PYZcb{}}\PY{l+s+s2}{. This corresponds to a wavelength of }\PY{l+s+si}{\PYZob{}}\PY{n}{wavelength}\PY{p}{(}\PY{n}{dE\PYZus{}covalent}\PY{p}{)}\PY{o}{.}\PY{n}{to}\PY{p}{(}\PY{n}{u}\PY{o}{.}\PY{n}{nm}\PY{p}{)}\PY{l+s+si}{:}\PY{l+s+s2}{.0f}\PY{l+s+si}{\PYZcb{}}\PY{l+s+s2}{,}\PY{l+s+se}{\PYZbs{}n}\PY{l+s+s2}{|      which is }\PY{l+s+si}{\PYZob{}}\PY{n}{spectrum\PYZus{}wavelength}\PY{p}{(}\PY{n}{wavelength}\PY{p}{(}\PY{n}{dE\PYZus{}covalent}\PY{p}{)}\PY{p}{)}\PY{l+s+si}{\PYZcb{}}\PY{l+s+s2}{. }\PY{l+s+se}{\PYZbs{}n}\PY{l+s+s2}{|}\PY{l+s+s2}{\PYZdq{}}\PY{p}{)}

\PY{n+nb}{print}\PY{p}{(}\PY{l+s+sa}{f}\PY{l+s+s2}{\PYZdq{}}\PY{l+s+s2}{|   b. The ionization energy of the hydrogen atom is }\PY{l+s+si}{\PYZob{}}\PY{n}{dE\PYZus{}H}\PY{l+s+si}{\PYZcb{}}\PY{l+s+s2}{ = }\PY{l+s+si}{\PYZob{}}\PY{n}{dE\PYZus{}H}\PY{o}{.}\PY{n}{to}\PY{p}{(}\PY{n}{u}\PY{o}{.}\PY{n}{J}\PY{p}{)}\PY{l+s+si}{:}\PY{l+s+s2}{.1e}\PY{l+s+si}{\PYZcb{}}\PY{l+s+s2}{.}\PY{l+s+se}{\PYZbs{}n}\PY{l+s+s2}{|      This corresponds to a wavelength of }\PY{l+s+si}{\PYZob{}}\PY{n}{wavelength}\PY{p}{(}\PY{n}{dE\PYZus{}H}\PY{p}{)}\PY{o}{.}\PY{n}{to}\PY{p}{(}\PY{n}{u}\PY{o}{.}\PY{n}{nm}\PY{p}{)}\PY{l+s+si}{:}\PY{l+s+s2}{.0f}\PY{l+s+si}{\PYZcb{}}\PY{l+s+s2}{, which is }\PY{l+s+si}{\PYZob{}}\PY{n}{spectrum\PYZus{}wavelength}\PY{p}{(}\PY{n}{wavelength}\PY{p}{(}\PY{n}{dE\PYZus{}H}\PY{p}{)}\PY{p}{)}\PY{l+s+si}{\PYZcb{}}\PY{l+s+s2}{. }\PY{l+s+se}{\PYZbs{}n}\PY{l+s+s2}{|}\PY{l+s+s2}{\PYZdq{}}\PY{p}{)}

\PY{n+nb}{print}\PY{p}{(}\PY{l+s+sa}{f}\PY{l+s+s2}{\PYZdq{}}\PY{l+s+s2}{|   c. This is the Lyman\PYZhy{}alpha line with energy }\PY{l+s+si}{\PYZob{}}\PY{n}{dE\PYZus{}12}\PY{l+s+si}{\PYZcb{}}\PY{l+s+s2}{ = }\PY{l+s+si}{\PYZob{}}\PY{n}{dE\PYZus{}12}\PY{o}{.}\PY{n}{to}\PY{p}{(}\PY{n}{u}\PY{o}{.}\PY{n}{J}\PY{p}{)}\PY{l+s+si}{:}\PY{l+s+s2}{.1e}\PY{l+s+si}{\PYZcb{}}\PY{l+s+s2}{.}\PY{l+s+se}{\PYZbs{}n}\PY{l+s+s2}{|      This corresponds to a wavelength of }\PY{l+s+si}{\PYZob{}}\PY{n}{wavelength}\PY{p}{(}\PY{n}{dE\PYZus{}12}\PY{p}{)}\PY{o}{.}\PY{n}{to}\PY{p}{(}\PY{n}{u}\PY{o}{.}\PY{n}{nm}\PY{p}{)}\PY{l+s+si}{:}\PY{l+s+s2}{.0f}\PY{l+s+si}{\PYZcb{}}\PY{l+s+s2}{, which is }\PY{l+s+si}{\PYZob{}}\PY{n}{spectrum\PYZus{}wavelength}\PY{p}{(}\PY{n}{wavelength}\PY{p}{(}\PY{n}{dE\PYZus{}12}\PY{p}{)}\PY{p}{)}\PY{l+s+si}{\PYZcb{}}\PY{l+s+s2}{. }\PY{l+s+se}{\PYZbs{}n}\PY{l+s+s2}{|}\PY{l+s+s2}{\PYZdq{}}\PY{p}{)}

\PY{n+nb}{print}\PY{p}{(}\PY{l+s+sa}{f}\PY{l+s+s2}{\PYZdq{}}\PY{l+s+s2}{|   d. The hydrogen\PYZhy{}alpha transition has energy }\PY{l+s+si}{\PYZob{}}\PY{n}{dE\PYZus{}23}\PY{l+s+si}{:}\PY{l+s+s2}{.1f}\PY{l+s+si}{\PYZcb{}}\PY{l+s+s2}{ = }\PY{l+s+si}{\PYZob{}}\PY{n}{dE\PYZus{}23}\PY{o}{.}\PY{n}{to}\PY{p}{(}\PY{n}{u}\PY{o}{.}\PY{n}{J}\PY{p}{)}\PY{l+s+si}{:}\PY{l+s+s2}{.1e}\PY{l+s+si}{\PYZcb{}}\PY{l+s+s2}{.}\PY{l+s+se}{\PYZbs{}n}\PY{l+s+s2}{|      This corresponds to a wavelength of }\PY{l+s+si}{\PYZob{}}\PY{n}{wavelength}\PY{p}{(}\PY{n}{dE\PYZus{}23}\PY{p}{)}\PY{o}{.}\PY{n}{to}\PY{p}{(}\PY{n}{u}\PY{o}{.}\PY{n}{nm}\PY{p}{)}\PY{l+s+si}{:}\PY{l+s+s2}{.0f}\PY{l+s+si}{\PYZcb{}}\PY{l+s+s2}{, which is }\PY{l+s+si}{\PYZob{}}\PY{n}{spectrum\PYZus{}wavelength}\PY{p}{(}\PY{n}{wavelength}\PY{p}{(}\PY{n}{dE\PYZus{}23}\PY{p}{)}\PY{p}{)}\PY{l+s+si}{\PYZcb{}}\PY{l+s+s2}{. }\PY{l+s+se}{\PYZbs{}n}\PY{l+s+s2}{|}\PY{l+s+s2}{\PYZdq{}}\PY{p}{)}

\PY{n+nb}{print}\PY{p}{(}\PY{l+s+sa}{f}\PY{l+s+s2}{\PYZdq{}}\PY{l+s+s2}{|   e.  See above.}\PY{l+s+se}{\PYZbs{}n}\PY{l+s+s2}{\PYZdq{}}\PY{p}{)}
\end{Verbatim}
\end{tcolorbox}

    \begin{Verbatim}[commandchars=\\\{\}]

|   ANSWER 4
|
|   a. Recall that the molecular hydrogen covalent bond has an energy of
|      4.5 eV = 7.2e-19 J. This corresponds to a wavelength of 276 nm,
|      which is ultraviolet.
|
|   b. The ionization energy of the hydrogen atom is 13.6 eV = 2.2e-18 J.
|      This corresponds to a wavelength of 91 nm, which is ultraviolet.
|
|   c. This is the Lyman-alpha line with energy 10.2 eV = 1.6e-18 J.
|      This corresponds to a wavelength of 122 nm, which is ultraviolet.
|
|   d. The hydrogen-alpha transition has energy 1.9 eV = 3.0e-19 J.
|      This corresponds to a wavelength of 656 nm, which is visible (red).
|
|   e.  See above.

    \end{Verbatim}

    Note that this calculation implies that it doesn't take much to break up
molecules.

The characterization of interstellar clouds can now be further refined.
Firstly, consider a dense molecular cloud. It contains dust grains,
molecules, and atoms all mixed together. Generally, hot young stars are
forming in these regions, and these stars emit copious amounts of
high-energy radiation. The higher energies do not penetrate into the
darkest parts of the cloud. We start deep in such a cloud and work our
way outward toward a star from there.

We start in regions with dust and molecular hydrogen (and other
molecules). Often the molecular hydrogen H2 will bind through van der
Waals forces to the dust, and the dust will shield it from dissociation
by the high-energy radiation. These regions can be quite cool at around
10 \texttt{K}. The dominant cooling mechanism is \textbf{molecular line
emission}: the ability of molecules to emit photons by changing their
vibrational and rotational states. Unfortunately, H2 is essentially
invisible, because it is so small and symmetrical: It is made of two
identical atoms so the center of mass is at the geometric center and
there is not asymmetry possible in its charge distribution, so it has no
electric dipole moment. This is important because to change from one
rotational state the next lower one would require a photon to carry away
one unit of angular momentum. However, photons can only do so by
coupling to the dipole moment, which H2 does not have.

What we need to look for is a different, less symmetrical molecule and
try to use that instead. Such a molecule is called a tracer, and they
are generally present in the ISM whenever conditions are favorable for
H2 formation. The most prevalent tracer for the dense, cool clouds we
are discussing here is CO. When it transitions from it's angular
momentum \(J=1\) state to the non-rotating \(J=0\) ground state, a
\(2.6\) \texttt{mm} photon is emitted.

 
\subsubsection{Problem 5}\label{problem-5}

\begin{enumerate}
\def\labelenumi{\alph{enumi}.}
\tightlist
\item
  Assuming this transition sets the scale for the temperature of the
  cloud, convert this photon energy to a temperature to confirm that it
  gives the correct order of magnitude.
\end{enumerate}
 

Continuing on our journey out of the molecular cloud toward a star, we
next encounter the \textbf{photodissociation region} (PDR): the region
where the penetrating far-UV radiation has sufficiently high energy to
dissociate H2 into H I atoms (problem 4a). The temperature in this
region is determined by the \ldots{} to be in the range 50 - 1000
\texttt{K} range. With the higher end closer to the star and the lower
end matching up with the dense molecular cloud, many layers are in the
100s of \texttt{K} range. While the far-UV radiation tends to heat the
gas, molecular line emission from C II and O I cool it. The wavelengths
of these modes are in the 10s-100s \textbf{microns}
(\(1\ \mu m = 10^{-6}\ m\)).

 
\subsubsection{Problem 5 (cont'd)}\label{problem-5-contd}

\begin{enumerate}
\def\labelenumi{\alph{enumi}.}
\setcounter{enumi}{1}
\tightlist
\item
  Use this to confirm that the claimed temperature of the PDR the
  correct order of magnitude.
\end{enumerate}
 

Finally, we get to regions in which the H I is exposed to the
high-energy radiation which ionizes it, or \textbf{photoionization}
region. This is the region where the dominant form of hydrogen is H II.
This region is characterized by a balance between the ejected electron
kinetic energy heating the gas, and the cooling by various ionic
transitions: The naïve conversion of the 91 \texttt{nm} wavelength of
the \emph{least} energetic \textbf{photoionizing} photon to a
temperature is misleading for two reasons. Firstly, the kinetic energy
of the ejected electron is \(hf+E_1=hf-13.6\) \texttt{eV} and not 13.6
\texttt{eV}. Secondly, there is again a different cooling mechanism,
this time predominantly from a forbidden O III transition line.

 
\subsubsection{Problem 5 (cont'd)}\label{problem-5-contd-1}

\begin{enumerate}
\def\labelenumi{\alph{enumi}.}
\setcounter{enumi}{2}
\tightlist
\item
  As mentioned in class, this line is green. Use this fact to estimate
  an upper limit for the temperature of the PDR.
\end{enumerate}
 

    \begin{tcolorbox}[breakable, size=fbox, boxrule=1pt, pad at break*=1mm,colback=cellbackground, colframe=cellborder]
\prompt{In}{incolor}{3}{\boxspacing}
\begin{Verbatim}[commandchars=\\\{\}]
\PY{k+kn}{from}\PY{+w}{ }\PY{n+nn}{astropy}\PY{n+nn}{.}\PY{n+nn}{constants}\PY{+w}{ }\PY{k+kn}{import} \PY{n}{k\PYZus{}B}

\PY{n+nb}{print}\PY{p}{(}\PY{l+s+sa}{f}\PY{l+s+s2}{\PYZdq{}}\PY{l+s+se}{\PYZbs{}n}\PY{l+s+s2}{|   ANSWER 5}\PY{l+s+se}{\PYZbs{}n}\PY{l+s+s2}{|}\PY{l+s+s2}{\PYZdq{}}\PY{p}{)}

\PY{k}{def}\PY{+w}{ }\PY{n+nf}{temperature}\PY{p}{(}\PY{n}{l}\PY{p}{:} \PY{n}{u}\PY{o}{.}\PY{n}{Quantity}\PY{p}{)} \PY{o}{\PYZhy{}}\PY{o}{\PYZgt{}} \PY{n}{u}\PY{o}{.}\PY{n}{Quantity}\PY{p}{:}
    \PY{n}{\PYZus{}t} \PY{o}{=} \PY{n}{h}\PY{o}{*}\PY{n}{c}\PY{o}{/}\PY{p}{(}\PY{n}{k\PYZus{}B}\PY{o}{*}\PY{n}{l}\PY{p}{)}
    \PY{k}{return} \PY{n}{\PYZus{}t}

\PY{n}{l\PYZus{}CO} \PY{o}{=} \PY{l+m+mf}{2.6}\PY{o}{*}\PY{n}{u}\PY{o}{.}\PY{n}{mm}
\PY{n}{l\PYZus{}micron} \PY{o}{=} \PY{l+m+mf}{100e\PYZhy{}6}\PY{o}{*}\PY{n}{u}\PY{o}{.}\PY{n}{m}
\PY{n}{l\PYZus{}green} \PY{o}{=} \PY{l+m+mi}{530}\PY{o}{*}\PY{n}{u}\PY{o}{.}\PY{n}{nm}


\PY{n+nb}{print}\PY{p}{(}\PY{l+s+sa}{f}\PY{l+s+s2}{\PYZdq{}}\PY{l+s+s2}{|   a. A wavelength of }\PY{l+s+si}{\PYZob{}}\PY{n}{l\PYZus{}CO}\PY{l+s+si}{\PYZcb{}}\PY{l+s+s2}{ corresponds to a temperature of }\PY{l+s+si}{\PYZob{}}\PY{n}{temperature}\PY{p}{(}\PY{n}{l\PYZus{}CO}\PY{p}{)}\PY{o}{.}\PY{n}{to}\PY{p}{(}\PY{n}{u}\PY{o}{.}\PY{n}{K}\PY{p}{)}\PY{l+s+si}{:}\PY{l+s+s2}{.1f}\PY{l+s+si}{\PYZcb{}}\PY{l+s+s2}{.}\PY{l+s+se}{\PYZbs{}n}\PY{l+s+s2}{|      This is indeed in the 10 K range.}\PY{l+s+se}{\PYZbs{}n}\PY{l+s+s2}{|}\PY{l+s+s2}{\PYZdq{}}\PY{p}{)}

\PY{n+nb}{print}\PY{p}{(}\PY{l+s+sa}{f}\PY{l+s+s2}{\PYZdq{}}\PY{l+s+s2}{|   b. A wavelength of 100 micron = }\PY{l+s+si}{\PYZob{}}\PY{n}{l\PYZus{}micron}\PY{l+s+si}{:}\PY{l+s+s2}{.1e}\PY{l+s+si}{\PYZcb{}}\PY{l+s+s2}{ corresponds to a temperature of }\PY{l+s+si}{\PYZob{}}\PY{n}{temperature}\PY{p}{(}\PY{n}{l\PYZus{}micron}\PY{p}{)}\PY{o}{.}\PY{n}{to}\PY{p}{(}\PY{n}{u}\PY{o}{.}\PY{n}{K}\PY{p}{)}\PY{l+s+si}{:}\PY{l+s+s2}{.0f}\PY{l+s+si}{\PYZcb{}}\PY{l+s+s2}{.}\PY{l+s+se}{\PYZbs{}n}\PY{l+s+s2}{|      This is indeed in the 100 K range.}\PY{l+s+se}{\PYZbs{}n}\PY{l+s+s2}{|}\PY{l+s+s2}{\PYZdq{}}\PY{p}{)}

\PY{n+nb}{print}\PY{p}{(}\PY{l+s+sa}{f}\PY{l+s+s2}{\PYZdq{}}\PY{l+s+s2}{|   c. Green is the 500\PYZhy{}565 nm range. Picking a wavelength of 530 nm gives a}\PY{l+s+se}{\PYZbs{}n}\PY{l+s+s2}{|      temperature of }\PY{l+s+si}{\PYZob{}}\PY{n}{temperature}\PY{p}{(}\PY{n}{l\PYZus{}green}\PY{p}{)}\PY{o}{.}\PY{n}{to}\PY{p}{(}\PY{n}{u}\PY{o}{.}\PY{n}{K}\PY{p}{)}\PY{l+s+si}{:}\PY{l+s+s2}{.1e}\PY{l+s+si}{\PYZcb{}}\PY{l+s+s2}{.}\PY{l+s+se}{\PYZbs{}n}\PY{l+s+s2}{\PYZdq{}}\PY{p}{)}
\end{Verbatim}
\end{tcolorbox}

    \begin{Verbatim}[commandchars=\\\{\}]

|   ANSWER 5
|
|   a. A wavelength of 2.6 mm corresponds to a temperature of 5.5 K.
|      This is indeed in the 10 K range.
|
|   b. A wavelength of 100 micron = 1.0e-04 m corresponds to a temperature of
144 K.
|      This is indeed in the 100 K range.
|
|   c. Green is the 500-565 nm range. Picking a wavelength of 530 nm gives a
|      temperature of 2.7e+04 K.

    \end{Verbatim}

    The last estimate is a bit high, but of the right dex. The typical H II
region temperatures usually quoted are in the range 7,000-10,000
\texttt{K}.

Above, we motivated the various regions by starting deep in a dense GMC
and then progressing toward a star so that successively more energetic
light results in various states of dissociation and ionization. However,
the atomic and molecular clouds described could also be considered as
astrophysical clouds in their own right. Below, we recapitulate our
considerations in problem 5, but as separate astrophysical phenomena.

\begin{itemize}
\tightlist
\item
  \href{https://en.wikipedia.org/wiki/Emission_nebula}{\textbf{Emission
  nebulae}} are clouds of ionized gas that are radiating the light
  absorbed from stars in the nebula. These are generally red due to the
  H\(\alpha\) line. There are two main types which are completely
  different:

  \begin{itemize}
  \tightlist
  \item
    \textbf{Stellar nurseries}: H II regions that are ionized from the
    UV radiation of new stars\\
  \item
    \href{https://en.wikipedia.org/wiki/Planetary_nebula}{\textbf{Planetary
    nebulae}}: Ionized gas remnants of dead/dying stars.{[}\^{}1{]}\\
  \end{itemize}
\item
  \href{https://en.wikipedia.org/wiki/Reflection_nebula}{\textbf{Reflection
  nebulae}}: Clouds of ISM that reflect light from nearby stars. Since
  short wavelength light is scattered more than long wavelength light,
  these nebulae tend to be bluish in color.
\item
  \href{https://en.wikipedia.org/wiki/Molecular_cloud}{\textbf{Atomic/molecular
  clouds}} come in two quite different types:

  \begin{itemize}
  \tightlist
  \item
    Molecular clouds' main hydrogen content is in the form of molecular
    hydrogen H2. They are usually dilute and cool \(\sim 10\)
    \texttt{K}.
  \item
    Atomic clouds are also called H I regions. Their temperatures can be
    in a wide range centered around an average of hundreds of Kelvin.
  \end{itemize}
\item
  \href{https://en.wikipedia.org/wiki/Dark_nebula}{\textbf{Absorption
  nebulae} or \textbf{dark nebulae}} are dust clouds that contain
  complex molecules like dust grains and are positioned in front of an
  optical source. Such clouds often look like holes in the sky because
  they scatter all light in the visible range away from our line of
  sight, although they tend to have reddish edges. Infrared and radio
  wave light may still pass through.
\end{itemize}

\subsubsection{Dipole interactions and the 21 cm
line}\label{dipole-interactions-and-the-21-cm-line}

We now come back to an important point, about H I regions that we
skipped over before. A hydrogen atom in its ground state consists of a
proton in the nucleus and an electron in the lowest orbital. As
mentioned in class, there are actually two states with closely-related
energies in this configuration: The spin of the electron and proton
could be aligned or anti-aligned. In \texttt{x\_spectroscopy}, we
described a spinning charged particle as a tiny magnet with a magnetic
dipole moment given by the Bohr magneton \[
\mu_B= \frac g2 \frac qm  s\hbar \tag{5}
\] Here \(q\) is the charge of the particle, \(m\) is its mass, and
\(s\) is its (half-integer) spin. The gyromagnetic ratio \(g\) is equal
to \(g=1\) ``classicaly'', but for a fundamental spin-1/2 (\(s=1/2\))
particle like the electron, \(g=2\). Although this has not mattered to
us before, the proton-although it too has \(s=1/2\)-is \emph{not} a
fundamental partcile in the same sense as the electron. Instead, it is a
composite of very strongly-interacting quarks and gluons in very complex
configurations. In fact the exact structure of the proton is an active
area of hadronic physics/strongly-coupled
\href{https://en.wikipedia.org/wiki/Quantum_chromodynamics}{Quantum
Chromodynamics (QCD)}. Experimentally, the gyromagnetic moment of the
proton is known to high accuracy: \(g_p = 5.5856946893(16)\) whereas
theoretically, lattice QCD and other techniques can only match this to
the few \% level.

As we discussed in \texttt{x\_spectroscopy}, the interaction energy of a
magnetic dipole \(\vec \mu\) and a magnetic field \(\vec B\) is given by
\(U=\vec \mu\cdot \vec B=\mu B \cos (\theta)\). (This ranges from
\(+\mu B\) (aligned) to \(-\mu B\) (anti-aligned), so the energy gap has
a size \(\Delta E=2\mu B\).) However, we have also mentioned that a
dipole acts like a little magnet, so it sets up its own magnetic field
(in fact a \emph{dipole field} by definition).

 
\subsubsection{Problem 6}\label{problem-6}

In this multi-part problem, we will explore dipoles and their
interactions, with the goal of understanding the 21 cm line. Start by
opening
\href{https://demo-clix.tiss.ac.in/phet/sims/html/charges-and-fields/latest/charges-and-fields_all.html?locale=en}{this}
simple on-line demo for electric charges and fields.

6a. We will make an electric dipole field by adding to monopole fields
(Coulomb charges) and letting the separation between the charges become
small.

\begin{enumerate}
\def\labelenumi{\roman{enumi}.}
\item
  Drag a \(+1\) \texttt{nC} charge (red) from the small box at the
  bottom and put it somewhere near the middle. Now remove it from the
  staging area (put it back in the storage box at the bottom), and
  repeat the exercise for the negative charge (blue). Describe the
  similarities and diffences in words. You do not have to be technical,
  but make your answers precise.
\item
  Starting over, drag a positive charge into the central area again and
  leave it there. Next, add a negative charge next to it to the left.
  Describe the shape of the resulting field as you change the
  separation.
\end{enumerate}

The little yellow dot labeled ``Sensors'' are what physicists call
``test charges''. A test charge is a fictitious charge that, when placed
in the electric field, feels a force in the direction of the field but
is assumed to be of such negible charge, that it does not contribute to
the electric field. (In reality it is not possible to sense a field
without affecting it, but you can limit the effect.)

\begin{enumerate}
\def\labelenumi{\roman{enumi}.}
\setcounter{enumi}{2}
\tightlist
\item
  After playing with a test charge in the dipole field, move it so that
  it the force it showes is parallel to the line connecting the charges.
  Leave the test charge there, note the size and direction of the force
  it feels, and then remove the negative charge. Explain how and why the
  force on the test charge changes. Put the negative charge back and
  then remove it again repeatedly until you have convinced yourself that
  the size of the force it feels is much smaller with the negative
  charge there than without it. What does this mean for the magnitude of
  the electric dipole field as a function of distance \(r\) from the
  dipole? Your answer will be qualitatitve, but try to make it precise.
  For example, how would you change the statement ``The magnitude of the
  monopole field falls off like \(1/r^2\).'' to one appropriate for the
  dipole field?
\end{enumerate}
 

     
\subsubsection{Answer 6}\label{answer-6}

 

    In the previous exercise, we modeled an \emph{electric} dipole field
using two electric monopoles (charges): one positive and one negative as
the two are brought close together. Although there are no magnetic
monopoles, the magnetic dipole field has the same mathematical form as
the electric one as long as the separation between monopoles is kept
very small.

Now, classically, the electron is a dipole in the magnetic field of the
proton or \emph{vice versa}. (From the definition of the dipole moment,
we expect the dipole moment of the proton to be much smaller than that
of the electron. (Why? Hint: look at the formula for the dipole moment
of a spinning point charge.) According to the formula for the energy of
a dipole in a magnetic field (of the other dipole), when the spins are
aligned, the energy of the system is slightly higher than when they are
anti-aligned.

\subsubsection{Problem 6 (cont'd)}\label{problem-6-contd}

In this problem, we will study the dipole-dipole interaction by looking
a the behavior of the field resulting from two aligned dipoles and that
resulting from two anti-aligned dipoles.

6b. Go back to the
\href{https://demo-clix.tiss.ac.in/phet/sims/html/charges-and-fields/latest/charges-and-fields_all.html?locale=en}{Phet
demo} and make a dipole out of two opposite charges as before.

\begin{enumerate}
\def\labelenumi{\roman{enumi}.}
\item
  Now put a second copy under it and parallel to it. Bring the two
  dipoles close together.
\item
  In a separate window, repeat exercise 6bi but with the \(\pm\) charges
  in the second copy switched.
\item
  Make a qualitative statement about what you notice about the differnce
  in the two resulting fields? In what sense is one bigger than the
  other?
\item
  Arrange 4 sensors symmetrically around your test charges in the same
  way in both fields. In what sense is the potential to move test
  charges around bigger in one field than that other.
\end{enumerate}

\subsubsection{Answer 6 (cont'd)}\label{answer-6-contd}


    The exercise above is supposed to convince you that when the spins of
the electron and proton are aligned, the energy of the system is
slightly higher than when they are anti-aligned, and we want to know the
wavelength of a photon emitted in a transition from the former to the
latter (or, equivalently absorbed by a tranition from the latter to the
former).

To recapitulate, a dipole has a magnetic field of its own. This field
falls off with distance \(r\) from the dipole faster than \(1/r^2\) as
we saw in problem 6aiii.

 
\subsubsection{Problem 6 (cont'd)}\label{problem-6-contd}

6c. Verify that the formula \[
{\vec B}(r) = x \frac {k}{c^2} \frac {\vec \mu}{r^3}
\] makes sense at least dimensionally. (\emph{I.e.} check that the
dimensions match on both sides.) Here \(x\) is just some numerical
coefficient we will determine later. Do this in the following steps

\begin{itemize}
\tightlist
\item
  Use the Lorentz force law in \texttt{x\_spectroscopy} to show that
  \([B]=\frac{M}{[e]T}\).
\item
  Remind yourself of \([k]\), for example from the Coulomb Law.
\item
  Calculate \(\left[ \frac {k}{c^2 r^3} \right]\).
\item
  Remind yourself of \([\mu]\) from \texttt{x\_spectroscopy} or above.
\item
  Put it all together.
\end{itemize}
 

     
\subsubsection{Answer 6 (cont'd)}\label{answer-6-contd}

6c. - From the Lorentz force law in \texttt{x\_spectroscopy}, \(F=qvB\)
allowing us to fix the dimensions of \(B\) as follows:
\(M\cdot L\cdot T^{-2} = [F]= [e] M\cdot T^{-1}[B]\). Solving for
\([B]\) we find the result claimed. - Next,
\(M\cdot L\cdot T^{-2} = [F] = [k]\cdot [e]^2 \cdot L^{-2}\), so
\([k]= [e]^{-2}\cdot M\cdot L^3\cdot T^{-2}\). - Calculating
\(\left[ \frac {k}{c^2 r^3} \right] = [e]^{-2}\cdot M\cdot L^{-2}\). -
From $\mu \sim \frac {e\hbar}m $, $[\mu]=[e]\cdot L^2 \cdot T^{-1}$.
- The RHS together gives
\([\textrm{RHS}]= [e]^{-1}\cdot M\cdot T^{-1}= [\textrm{LHS}]\), as
claimed.
 

    From the formula above, we get the magnetic field of the proton, and
therefore the (magnetic contribution to the) energy of an electron in
the proton field. Evaluating this at the Bohr radius gives an additional
contribution to the energy of \[
x \frac{k \vec \mu_e \cdot \vec \mu_p}{c^2 a^3} 
%\frac 83 \frac{k}{c^2}\frac{\vec \mu_e \cdot \vec \mu_p}{a^3} 
\] This form is pretty nice, but we can make it even more illuminating.

 
\subsubsection{Problem 6 (cont'd)}\label{problem-6-contd}

6d. Define the energy \(E_\textrm{hf}= x \frac{k\mu_e \mu_p}{c^2 a^3}\)
and plug in the formulæ for the magnetic moments to derive this energy
in the following forms:

\begin{enumerate}
\def\labelenumi{\roman{enumi}.}
\item
  First show you can write it as \[
  x \frac{ke^2}{a}\frac1{c^2 a^2} \frac{g_p}{8} \frac{\hbar^2}{m_p m_e} 
  = \frac{x g_p}{8}\frac{ke^2}{a}\frac{m_e}{m_p}\left( \frac {\hbar}{m_e c a} \right)^2
  \]
\item
  From \texttt{x\_spectroscopy} remind yourself of the simple formulæ we
  derived for ionization energy of hydrogen \(E_1\) and the find
  structure constant \(\alpha\), and use them to show that \[
  E_\textrm{hf}
  = \frac{x g_p}{8} \frac{m_e}{m_p} \left|E_1\right| \alpha^2 
  \]
\item
  Recall that there were actually two configuations with energies
  \(+\left|E_1\right|\) (aligned) and \(-\left|E_1\right|\)
  (anti-aligned), and use it to conclude that in a transition from the
  former to the latter, the energy of the emitted photon must be equal
  to \[
  \Delta E_\textrm{hf}
  = \frac{x g_p}{4} \frac{m_e}{m_p} \left|E_1\right| \alpha^2   \tag{6}
  \]
\end{enumerate}
 

     
\subsubsection{Answer 6 (cont'd)}\label{answer-6-contd}

6di. Plugging in \(\mu = \frac{g_pe\hbar}{4m_p}\) and
\(\mu_e = \frac{e\hbar}{2m_e}\) into the formula gives
\[E_\textrm{hf}= x \frac{k g_p e^2\hbar^2}{8 m_p m_e c^2 a^3}\] We can
separate this as indicated in the question, but why would we? The answer
is that it clarifies some things: * the factor \(\frac{ke^2}{a}\) is
familiar from as the potential energy of the Bohr atom in the ground
state, so the rest should be dimensionless.

\begin{itemize}
\item
  The factor \(\frac{g_p}{8}\) is a dimensionless number, so that leaves
  \(\frac1{c^2 a^2}\frac{\hbar^2}{m_p m_e}\) which can be separated
  democratically into
\item
  a factor \(\frac{\hbar}{m_eca}\) for the electron and
\item
  another \(\frac{\hbar}{m_pca}\) for the proton.
\end{itemize}

Are these dimensionless? You bet: the \(m_e c\) factor is a momentum and
\([\hbar] = [\textrm{momentum}]\cdot L\), so this is dimensionless.
Bells should be ringing loudly now, because we have been here before:
This is the fine-structure constant \(\alpha\)!

\begin{enumerate}
\def\labelenumi{\roman{enumi}.}
\setcounter{enumi}{1}
\item
  We already covered this above: the factor \(\frac{ke^2}{a} = |E_1|\).
  That leaves the factor \(\frac{\hbar}{m_pca}\). If we rewrite this as
  \(\frac{\hbar}{m_pca}=\frac{m_e}{m_p}\frac{\hbar}{m_eca}=\frac{m_e}{m_p}\alpha\),
  then we get another factor of \(\alpha\). There is also a factor of
  the ratio of the masses, which is a small number
  \(\sim \frac 1{1800}\). Thus, we have found the requested form.
\item
  The two energies for the aligned state (\(\theta=0\)) and anti-aligned
  state (\(\theta=\pi\)) which differ by a factor of
  \(\cos(0)-\cos(\pi)=2\), so the result follows.
\end{enumerate}
 

     
\subsubsection{Problem 6 (cont'd)}\label{problem-6-contd}

6e. A course on classical electrodynamics suffices to fix the
coefficient \(x=\frac{32}3\).

\begin{enumerate}
\def\labelenumi{\roman{enumi}.}
\item
  Plug in the numbers to show that the energy gap is
  \(\Delta E_{hf} = 5.88\times 10^{-6}\) \texttt{eV}
\item
  Convert this energy to a wavelength in centimeters. What part of the
  spectrum is this?
\end{enumerate}
 

     
\subsubsection{Answer 6 (cont'd)}\label{answer-6-contd}

6ei. We can do this various ways. We choose to define plug in the final
form (6) derived in problem 6d.
 

    \begin{tcolorbox}[breakable, size=fbox, boxrule=1pt, pad at break*=1mm,colback=cellbackground, colframe=cellborder]
\prompt{In}{incolor}{4}{\boxspacing}
\begin{Verbatim}[commandchars=\\\{\}]
\PY{k+kn}{from}\PY{+w}{ }\PY{n+nn}{astropy}\PY{n+nn}{.}\PY{n+nn}{constants}\PY{+w}{ }\PY{k+kn}{import} \PY{n}{hbar}\PY{p}{,} \PY{n}{eps0}\PY{p}{,} \PY{n}{e}\PY{p}{,} \PY{n}{m\PYZus{}e}\PY{p}{,} \PY{n}{m\PYZus{}p}
\PY{k+kn}{from}\PY{+w}{ }\PY{n+nn}{math}\PY{+w}{ }\PY{k+kn}{import} \PY{n}{pi}

\PY{n}{e} \PY{o}{=} \PY{n}{e}\PY{o}{.}\PY{n}{to}\PY{p}{(}\PY{n}{u}\PY{o}{.}\PY{n}{C}\PY{p}{)}
\PY{n}{k} \PY{o}{=} \PY{l+m+mi}{1}\PY{o}{/}\PY{p}{(}\PY{l+m+mi}{4}\PY{o}{*}\PY{n}{pi}\PY{o}{*}\PY{n}{eps0}\PY{p}{)}
\PY{n}{alpha} \PY{o}{=}\PY{p}{(} \PY{p}{(}\PY{n}{k}\PY{o}{*}\PY{n}{e}\PY{o}{*}\PY{o}{*}\PY{l+m+mi}{2}\PY{p}{)}\PY{o}{/}\PY{p}{(}\PY{n}{hbar}\PY{o}{*}\PY{n}{c}\PY{p}{)} \PY{p}{)}\PY{o}{.}\PY{n}{decompose}\PY{p}{(}\PY{p}{)}
\PY{n}{E\PYZus{}1} \PY{o}{=} \PY{l+m+mi}{1}\PY{o}{/}\PY{l+m+mi}{2}\PY{o}{*}\PY{n}{m\PYZus{}e}\PY{o}{*}\PY{p}{(}\PY{n}{alpha}\PY{o}{*}\PY{n}{c}\PY{p}{)}\PY{o}{*}\PY{o}{*}\PY{l+m+mi}{2}

\PY{k}{def}\PY{+w}{ }\PY{n+nf}{E\PYZus{}hf}\PY{p}{(}\PY{n}{x}\PY{p}{:} \PY{n+nb}{float}\PY{p}{,} \PY{n}{g}\PY{p}{:} \PY{n+nb}{float}\PY{p}{)}\PY{o}{\PYZhy{}}\PY{o}{\PYZgt{}} \PY{n}{u}\PY{o}{.}\PY{n}{Quantity}\PY{p}{:}
    \PY{n}{\PYZus{}E} \PY{o}{=} \PY{l+m+mf}{0.25}\PY{o}{*}\PY{n}{x}\PY{o}{*}\PY{n}{g}\PY{o}{*}\PY{p}{(}\PY{n}{m\PYZus{}e}\PY{o}{/}\PY{n}{m\PYZus{}p}\PY{p}{)}\PY{o}{*}\PY{n}{E\PYZus{}1}\PY{o}{*}\PY{n}{alpha}\PY{o}{*}\PY{o}{*}\PY{l+m+mi}{2}
    \PY{k}{return} \PY{n}{\PYZus{}E}\PY{o}{.}\PY{n}{to}\PY{p}{(}\PY{n}{u}\PY{o}{.}\PY{n}{J}\PY{p}{)}

\PY{n+nb}{print}\PY{p}{(}\PY{l+s+sa}{f}\PY{l+s+s2}{\PYZdq{}}\PY{l+s+se}{\PYZbs{}n}\PY{l+s+s2}{|   ANSWER 6 (cont}\PY{l+s+s2}{\PYZsq{}}\PY{l+s+s2}{d)}\PY{l+s+se}{\PYZbs{}n}\PY{l+s+s2}{|}\PY{l+s+s2}{\PYZdq{}}\PY{p}{)}

\PY{n}{\PYZus{}E21} \PY{o}{=} \PY{n}{E\PYZus{}hf}\PY{p}{(} \PY{l+m+mi}{32}\PY{o}{/}\PY{l+m+mi}{3}\PY{p}{,} \PY{l+m+mf}{5.59} \PY{p}{)}

\PY{n+nb}{print}\PY{p}{(}\PY{l+s+sa}{f}\PY{l+s+s2}{\PYZdq{}}\PY{l+s+s2}{|   i. With x=32/3 and g\PYZus{}p=5.59 for the proton, as reported above, the gap energy is}\PY{l+s+se}{\PYZbs{}n}\PY{l+s+s2}{|   ΔE\PYZus{}hf = }\PY{l+s+si}{\PYZob{}}\PY{n}{\PYZus{}E21}\PY{o}{.}\PY{n}{to}\PY{p}{(}\PY{n}{u}\PY{o}{.}\PY{n}{eV}\PY{p}{)}\PY{l+s+si}{:}\PY{l+s+s2}{.2e}\PY{l+s+si}{\PYZcb{}}\PY{l+s+s2}{,}\PY{l+s+se}{\PYZbs{}n}\PY{l+s+s2}{|}\PY{l+s+s2}{\PYZdq{}}\PY{p}{)}

\PY{n+nb}{print}\PY{p}{(}\PY{l+s+sa}{f}\PY{l+s+s2}{\PYZdq{}}\PY{l+s+s2}{|   ii. The corresponding wavelength is }\PY{l+s+si}{\PYZob{}}\PY{n}{wavelength}\PY{p}{(}\PY{n}{\PYZus{}E21}\PY{p}{)}\PY{o}{.}\PY{n}{to}\PY{p}{(}\PY{n}{u}\PY{o}{.}\PY{n}{cm}\PY{p}{)}\PY{l+s+si}{:}\PY{l+s+s2}{.0f}\PY{l+s+si}{\PYZcb{}}\PY{l+s+s2}{.}\PY{l+s+se}{\PYZbs{}n}\PY{l+s+s2}{|   This is a }\PY{l+s+si}{\PYZob{}}\PY{n}{spectrum\PYZus{}wavelength}\PY{p}{(}\PY{n}{wavelength}\PY{p}{(}\PY{n}{\PYZus{}E21}\PY{p}{)}\PY{p}{)}\PY{l+s+si}{\PYZcb{}}\PY{l+s+s2}{. }\PY{l+s+se}{\PYZbs{}n}\PY{l+s+s2}{\PYZdq{}}\PY{p}{)}
\end{Verbatim}
\end{tcolorbox}

    \begin{Verbatim}[commandchars=\\\{\}]

|   ANSWER 6 (cont'd)
|
|   i. With x=32/3 and g\_p=5.59 for the proton, as reported above, the gap
energy is
|   ΔE\_hf = 5.88e-06 eV,
|
|   ii. The corresponding wavelength is 21 cm.
|   This is a radio wave.

    \end{Verbatim}

     
\subsubsection{Answer 6 (cont'd)}\label{answer-6-contd}

6ei. Alternatively, we can do it from the original dipole-dipole
interaction.
 

    \begin{tcolorbox}[breakable, size=fbox, boxrule=1pt, pad at break*=1mm,colback=cellbackground, colframe=cellborder]
\prompt{In}{incolor}{5}{\boxspacing}
\begin{Verbatim}[commandchars=\\\{\}]
\PY{n}{a} \PY{o}{=} \PY{p}{(} \PY{n}{hbar}\PY{o}{/}\PY{p}{(}\PY{n}{m\PYZus{}e}\PY{o}{*}\PY{n}{alpha}\PY{o}{*}\PY{n}{c}\PY{p}{)} \PY{p}{)}\PY{o}{.}\PY{n}{decompose}\PY{p}{(}\PY{p}{)}

\PY{k}{def}\PY{+w}{ }\PY{n+nf}{mu}\PY{p}{(}\PY{n}{g}\PY{p}{,}\PY{n}{q}\PY{p}{,}\PY{n}{m}\PY{p}{:} \PY{n}{u}\PY{o}{.}\PY{n}{Quantity}\PY{p}{,} \PY{n}{s}\PY{p}{:} \PY{n+nb}{float}\PY{p}{)}\PY{o}{\PYZhy{}}\PY{o}{\PYZgt{}} \PY{n}{u}\PY{o}{.}\PY{n}{Quantity}\PY{p}{:}
    \PY{n}{\PYZus{}mu} \PY{o}{=} \PY{p}{(}\PY{n}{g}\PY{o}{/}\PY{l+m+mi}{2}\PY{p}{)}\PY{o}{*}\PY{p}{(}\PY{n}{q}\PY{o}{/}\PY{n}{m}\PY{p}{)}\PY{o}{*}\PY{n}{s}\PY{o}{*}\PY{n}{hbar} 
    \PY{k}{return} \PY{n}{\PYZus{}mu}\PY{o}{.}\PY{n}{decompose}\PY{p}{(}\PY{p}{)}

\PY{n}{mu\PYZus{}e} \PY{o}{=} \PY{n}{mu}\PY{p}{(}\PY{l+m+mi}{2}\PY{p}{,}\PY{n}{e}\PY{p}{,}\PY{n}{m\PYZus{}e}\PY{p}{,}\PY{l+m+mi}{1}\PY{o}{/}\PY{l+m+mi}{2}\PY{p}{)}
\PY{n}{mu\PYZus{}p} \PY{o}{=} \PY{n}{mu}\PY{p}{(}\PY{l+m+mf}{5.59}\PY{p}{,}\PY{n}{e}\PY{p}{,}\PY{n}{m\PYZus{}p}\PY{p}{,}\PY{l+m+mi}{1}\PY{o}{/}\PY{l+m+mi}{2}\PY{p}{)}

\PY{k}{def}\PY{+w}{ }\PY{n+nf}{E\PYZus{}dipole}\PY{p}{(}\PY{n}{mu\PYZus{}1}\PY{p}{:} \PY{n}{u}\PY{o}{.}\PY{n}{Quantity}\PY{p}{,} \PY{n}{mu\PYZus{}2}\PY{p}{:} \PY{n}{u}\PY{o}{.}\PY{n}{Quantity}\PY{p}{,} \PY{n}{r}\PY{p}{:} \PY{n}{u}\PY{o}{.}\PY{n}{Quantity}\PY{p}{)}\PY{o}{\PYZhy{}}\PY{o}{\PYZgt{}} \PY{n}{u}\PY{o}{.}\PY{n}{Quantity}\PY{p}{:}
    \PY{n}{\PYZus{}E} \PY{o}{=} \PY{l+m+mi}{32}\PY{o}{/}\PY{l+m+mi}{3}\PY{o}{*}\PY{n}{k}\PY{o}{/}\PY{p}{(}\PY{n}{c}\PY{o}{*}\PY{o}{*}\PY{l+m+mi}{2}\PY{p}{)}\PY{o}{*}\PY{n}{mu\PYZus{}1}\PY{o}{*}\PY{n}{mu\PYZus{}2}\PY{o}{/}\PY{n}{r}\PY{o}{*}\PY{o}{*}\PY{l+m+mi}{3}
    \PY{k}{return} \PY{n}{\PYZus{}E}\PY{o}{.}\PY{n}{to}\PY{p}{(}\PY{n}{u}\PY{o}{.}\PY{n}{J}\PY{p}{)}

\PY{k}{def}\PY{+w}{ }\PY{n+nf}{E\PYZus{}hf}\PY{p}{(}\PY{n}{x}\PY{p}{:} \PY{n+nb}{float}\PY{p}{,} \PY{n}{g}\PY{p}{:} \PY{n+nb}{float}\PY{p}{)}\PY{o}{\PYZhy{}}\PY{o}{\PYZgt{}} \PY{n}{u}\PY{o}{.}\PY{n}{Quantity}\PY{p}{:}
    \PY{n}{\PYZus{}E} \PY{o}{=} \PY{l+m+mf}{0.25}\PY{o}{*}\PY{n}{x}\PY{o}{*}\PY{n}{g}\PY{o}{*}\PY{p}{(}\PY{n}{m\PYZus{}e}\PY{o}{/}\PY{n}{m\PYZus{}p}\PY{p}{)}\PY{o}{*}\PY{n}{E\PYZus{}1}\PY{o}{*}\PY{n}{alpha}\PY{o}{*}\PY{o}{*}\PY{l+m+mi}{2}
    \PY{k}{return} \PY{n}{\PYZus{}E}\PY{o}{.}\PY{n}{to}\PY{p}{(}\PY{n}{u}\PY{o}{.}\PY{n}{J}\PY{p}{)}

\PY{n+nb}{print}\PY{p}{(}\PY{l+s+sa}{f}\PY{l+s+s2}{\PYZdq{}}\PY{l+s+se}{\PYZbs{}n}\PY{l+s+s2}{|   ANSWER 6 (cont}\PY{l+s+s2}{\PYZsq{}}\PY{l+s+s2}{d)}\PY{l+s+se}{\PYZbs{}n}\PY{l+s+s2}{|}\PY{l+s+s2}{\PYZdq{}}\PY{p}{)}

\PY{n+nb}{print}\PY{p}{(}\PY{l+s+sa}{f}\PY{l+s+s2}{\PYZdq{}}\PY{l+s+s2}{|   E\PYZus{}dipole = }\PY{l+s+si}{\PYZob{}}\PY{n}{E\PYZus{}dipole}\PY{p}{(}\PY{n}{mu\PYZus{}e}\PY{p}{,}\PY{+w}{ }\PY{n}{mu\PYZus{}p}\PY{p}{,}\PY{+w}{ }\PY{n}{a}\PY{p}{)}\PY{o}{.}\PY{n}{to}\PY{p}{(}\PY{n}{u}\PY{o}{.}\PY{n}{eV}\PY{p}{)}\PY{l+s+si}{:}\PY{l+s+s2}{.2e}\PY{l+s+si}{\PYZcb{}}\PY{l+s+s2}{.}\PY{l+s+se}{\PYZbs{}n}\PY{l+s+s2}{\PYZdq{}}\PY{p}{)}
\end{Verbatim}
\end{tcolorbox}

    \begin{Verbatim}[commandchars=\\\{\}]

|   ANSWER 6 (cont'd)
|
|   E\_dipole = 5.88e-06 eV.

    \end{Verbatim}

    This is called the
\href{https://en.wikipedia.org/wiki/Hydrogen_line}{\textbf{21-cm line}}
or H I line. It is part of the so-called hyperfine structure of the
hydrogen atom: The exceedingly closely-spaced energy levels between
states in the hydrogen spectrum. The energy difference is so small that
the spins can stay aligned for \emph{millions} of years. This is an
incredibly long time especially when you consider that particles in
gasses typically collide \emph{billions} of times every second. These
conditions disrupt the excited state, implying that we do not see
certain transitions in laboratory experiments. In typically dramatic
fashion, we call these processes
``\href{https://en.wikipedia.org/wiki/Forbidden_mechanism\#In_astrophysics_and_atomic_physics}{\textbf{forbidden}}''.
In the exceedingly rarified gasses encountered in space, the collision
time is measured in \emph{years}, not nanoseconds. Statistically, then,
if we have enough particles some of them will have time to decay,
thereby emitting a photon with a wavelength of 21-cm.

 
\subsubsection{Problem 7}\label{problem-7}

Using table 1, estimate the order of magnitude of the number of
particles in an interstellar cloud.
 

    \begin{tcolorbox}[breakable, size=fbox, boxrule=1pt, pad at break*=1mm,colback=cellbackground, colframe=cellborder]
\prompt{In}{incolor}{6}{\boxspacing}
\begin{Verbatim}[commandchars=\\\{\}]
\PY{k+kn}{from}\PY{+w}{ }\PY{n+nn}{astropy}\PY{n+nn}{.}\PY{n+nn}{constants}\PY{+w}{ }\PY{k+kn}{import} \PY{n}{M\PYZus{}sun}\PY{p}{,} \PY{n}{m\PYZus{}p}

\PY{n+nb}{print}\PY{p}{(}\PY{l+s+sa}{f}\PY{l+s+s2}{\PYZdq{}}\PY{l+s+se}{\PYZbs{}n}\PY{l+s+s2}{|   ANSWER 7}\PY{l+s+se}{\PYZbs{}n}\PY{l+s+s2}{|}\PY{l+s+s2}{\PYZdq{}}\PY{p}{)}

\PY{n}{\PYZus{}d} \PY{o}{=} \PY{l+m+mf}{1e17}\PY{o}{*}\PY{n}{u}\PY{o}{.}\PY{n}{m}
\PY{n}{\PYZus{}V} \PY{o}{=} \PY{n}{\PYZus{}d}\PY{o}{*}\PY{o}{*}\PY{l+m+mi}{3}
\PY{n}{\PYZus{}rho} \PY{o}{=} \PY{l+m+mf}{1e\PYZhy{}18}\PY{o}{*}\PY{n}{u}\PY{o}{.}\PY{n}{kg}\PY{o}{/}\PY{n}{u}\PY{o}{.}\PY{n}{m}\PY{o}{*}\PY{o}{*}\PY{l+m+mi}{3}
\PY{n}{\PYZus{}m} \PY{o}{=} \PY{n}{\PYZus{}rho}\PY{o}{*}\PY{n}{\PYZus{}V}

\PY{n+nb}{print}\PY{p}{(}\PY{l+s+sa}{f}\PY{l+s+s2}{\PYZdq{}}\PY{l+s+s2}{|   At a size of d \PYZti{} }\PY{l+s+si}{\PYZob{}}\PY{n}{\PYZus{}d}\PY{l+s+si}{:}\PY{l+s+s2}{.0e}\PY{l+s+si}{\PYZcb{}}\PY{l+s+s2}{, the volume is of order }\PY{l+s+si}{\PYZob{}}\PY{n}{\PYZus{}V}\PY{l+s+si}{:}\PY{l+s+s2}{.0e}\PY{l+s+si}{\PYZcb{}}\PY{l+s+s2}{. At a density}\PY{l+s+se}{\PYZbs{}n}\PY{l+s+s2}{|   rho \PYZti{} }\PY{l+s+si}{\PYZob{}}\PY{n}{\PYZus{}rho}\PY{l+s+si}{:}\PY{l+s+s2}{.0e}\PY{l+s+si}{\PYZcb{}}\PY{l+s+s2}{, this is a mass of }\PY{l+s+si}{\PYZob{}}\PY{n}{\PYZus{}m}\PY{l+s+si}{:}\PY{l+s+s2}{.0e}\PY{l+s+si}{\PYZcb{}}\PY{l+s+s2}{ or \PYZti{} }\PY{l+s+si}{\PYZob{}}\PY{n+nb}{int}\PY{p}{(}\PY{n+nb}{round}\PY{p}{(}\PY{n}{\PYZus{}m}\PY{o}{/}\PY{n}{M\PYZus{}sun}\PY{p}{,}\PY{o}{\PYZhy{}}\PY{l+m+mi}{1}\PY{p}{)}\PY{p}{)}\PY{l+s+si}{\PYZcb{}}\PY{l+s+s2}{ solar masses.}\PY{l+s+se}{\PYZbs{}n}\PY{l+s+s2}{|   About 3/4 is protons and 1/4 is alphas (by mass), but for this back\PYZhy{}of\PYZhy{}the\PYZhy{}}\PY{l+s+se}{\PYZbs{}n}\PY{l+s+s2}{|   envelope calculation, this doesn}\PY{l+s+s2}{\PYZsq{}}\PY{l+s+s2}{t matter. Dividing by the mass m\PYZus{}p = }\PY{l+s+si}{\PYZob{}}\PY{n}{m\PYZus{}p}\PY{l+s+si}{:}\PY{l+s+s2}{.0e}\PY{l+s+si}{\PYZcb{}}\PY{l+s+se}{\PYZbs{}n}\PY{l+s+s2}{|   of the proton, we find that such a cloud contains on the order of N \PYZti{} }\PY{l+s+si}{\PYZob{}}\PY{n}{\PYZus{}m}\PY{o}{/}\PY{n}{m\PYZus{}p}\PY{l+s+si}{:}\PY{l+s+s2}{.0e}\PY{l+s+si}{\PYZcb{}}\PY{l+s+s2}{ or}\PY{l+s+se}{\PYZbs{}n}\PY{l+s+s2}{|   10\PYZca{}60 hydrogen atoms.}\PY{l+s+se}{\PYZbs{}n}\PY{l+s+s2}{\PYZdq{}}\PY{p}{)}
\end{Verbatim}
\end{tcolorbox}

    \begin{Verbatim}[commandchars=\\\{\}]

|   ANSWER 7
|
|   At a size of d \textasciitilde{} 1e+17 m, the volume is of order 1e+51 m3. At a density
|   rho \textasciitilde{} 1e-18 kg / m3, this is a mass of 1e+33 kg or \textasciitilde{} 500 solar masses.
|   About 3/4 is protons and 1/4 is alphas (by mass), but for this back-of-the-
|   envelope calculation, this doesn't matter. Dividing by the mass m\_p = 2e-27
kg
|   of the proton, we find that such a cloud contains on the order of N \textasciitilde{} 6e+59
or
|   10\^{}60 hydrogen atoms.

    \end{Verbatim}

    A statistical calculation is necessary to find the rate of emission of
21-cm radiation from an interstellar cloud, but this number of particles
is so large that we can expect a measurable rate. The 21-cm line was
predicted to be measurable in 1944, and was detected in 1951. It is the
hallmark of ISM, and was used in 1952 to map out the distribution of
matter in the galaxy. Not only can the presence of H I be detected, but
the Doppler shift can be used to find radial velocity. Even the
interstellar magnetic fields can be determined, because magnetic fields
change the energy levels of the hydrogen atom, and therefore further
shift the emission spectrum (remember the
\href{https://en.wikipedia.org/wiki/Zeeman_effect}{Zeeman effect}? psst:
the final is cumulative\ldots{}).

The figure above-right shows the
\href{https://en.wikipedia.org/wiki/Pioneer_plaque}{Pioneer plaque} and
the one below right shows the
\href{https://en.wikipedia.org/wiki/Voyager_Golden_Record}{Voyager
Golden Record}. The diagrams at the top left and bottom right
respectively are supposed to represent the hyperfine transition of
hydrogen graphically. These objects were attached to the probes in case
those probes were ever intercepted by an extraterrestrial intelligence.

\begin{center}\rule{0.5\linewidth}{0.5pt}\end{center}

\subsection{Stellar formation}\label{stellar-formation}

The birth of a star proceeds in the seven stages summarized in table 1.
In the first three stages, an interstellar cloud of atoms, dust and/or
molecules contracts and fragments with a fragment further contracting
into the initial stage of a protostar.

\begin{longtable}[]{@{}
  >{\centering\arraybackslash}p{(\linewidth - 14\tabcolsep) * \real{0.1250}}
  >{\centering\arraybackslash}p{(\linewidth - 14\tabcolsep) * \real{0.1250}}
  >{\centering\arraybackslash}p{(\linewidth - 14\tabcolsep) * \real{0.1250}}
  >{\centering\arraybackslash}p{(\linewidth - 14\tabcolsep) * \real{0.1250}}
  >{\centering\arraybackslash}p{(\linewidth - 14\tabcolsep) * \real{0.1250}}
  >{\centering\arraybackslash}p{(\linewidth - 14\tabcolsep) * \real{0.1250}}
  >{\centering\arraybackslash}p{(\linewidth - 14\tabcolsep) * \real{0.1250}}
  >{\centering\arraybackslash}p{(\linewidth - 14\tabcolsep) * \real{0.1250}}@{}}
\toprule\noalign{}
\begin{minipage}[b]{\linewidth}\centering
Stage
\end{minipage} & \begin{minipage}[b]{\linewidth}\centering
time (yr)
\end{minipage} & \begin{minipage}[b]{\linewidth}\centering
\(T_\textrm{core}\)(K)
\end{minipage} & \begin{minipage}[b]{\linewidth}\centering
\(T_\textrm{surface}\) (K)
\end{minipage} & \begin{minipage}[b]{\linewidth}\centering
\(\rho_\textrm{core}\) g/cm3
\end{minipage} & \begin{minipage}[b]{\linewidth}\centering
diameter (km)
\end{minipage} & \begin{minipage}[b]{\linewidth}\centering
comparable object
\end{minipage} & \begin{minipage}[b]{\linewidth}\centering
object
\end{minipage} \\
\midrule\noalign{}
\endhead
\bottomrule\noalign{}
\endlastfoot
1 & \(10^6\) & 10 & 10 & \(1.7\times 10^{-21}\) & \(10^{14}\sim 3\) pc &
nearest stars & interstellar cloud \\
2 & \(10^4\) & 100 & 10 & \(1.7\times 10^{-18}\) & \(10^{12}\) & --- &
cloud fragment \\
3 & \(10^5\) & 10,000 & 100 & \(1.7\times 10^{-12}\) &
\(10^{10}\sim 100\) au & solar system & cloud frag/ protostar \\
4 & \(10^6\) & 1,000,000 & 3000 & \(1.7\times 10^{-6}\) &
\(10^8\sim 100 D_\odot\) & earth orbit & PMS star \\
5 & \(10^7\) & 5,000,000 & 4000 & \(0.017\) & \(10^7 \sim 10 D_\odot\) &
--- & PMS star \\
6 & \(10^7\) & 10,000,000 & 4500 & \(17\) & \(D_\odot\) & sun & star \\
7 & \(10^{10}\) & 15,000,000 & 6000 & \(170\) & \(D_\odot\) & sun & star
(main sequence) \\
\end{longtable}

Table 1: Stages of stellar formation. For reference
\(D_\odot=1.4\times 10^6\) km, 1 pc \(=3.1\times 10^{13}\) km, and the
diameter of the Kuiper belt 100 AU \(=1.5\times 10^{10}\) km. The core
density of the sun is \(\rho_{\textrm{core},\odot} = 160\) g/cm3. PMS
stands for pre-main sequence.

\subsubsection{Stages 1-3}\label{stages-1-3}

We will be brief about these stages. Derivations of stated results can
be found in \texttt{ss\_solar\_system\_1,2,3}. We will start with the
ISM context and progressing to smaller subregions down to that dense
clump of the \textbf{giant molecular cloud} (GMC) which eventually
collapses to a protostar. The structures may be represented by the
following table.

\begin{longtable}[]{@{}
  >{\raggedright\arraybackslash}p{(\linewidth - 18\tabcolsep) * \real{0.0882}}
  >{\raggedright\arraybackslash}p{(\linewidth - 18\tabcolsep) * \real{0.0882}}
  >{\centering\arraybackslash}p{(\linewidth - 18\tabcolsep) * \real{0.1176}}
  >{\raggedright\arraybackslash}p{(\linewidth - 18\tabcolsep) * \real{0.0882}}
  >{\centering\arraybackslash}p{(\linewidth - 18\tabcolsep) * \real{0.1176}}
  >{\raggedright\arraybackslash}p{(\linewidth - 18\tabcolsep) * \real{0.0882}}
  >{\centering\arraybackslash}p{(\linewidth - 18\tabcolsep) * \real{0.1176}}
  >{\raggedright\arraybackslash}p{(\linewidth - 18\tabcolsep) * \real{0.0882}}
  >{\centering\arraybackslash}p{(\linewidth - 18\tabcolsep) * \real{0.1176}}
  >{\raggedright\arraybackslash}p{(\linewidth - 18\tabcolsep) * \real{0.0882}}@{}}
\toprule\noalign{}
\begin{minipage}[b]{\linewidth}\raggedright
\end{minipage} & \begin{minipage}[b]{\linewidth}\raggedright
ISM complex
\end{minipage} & \begin{minipage}[b]{\linewidth}\centering
\(\supset\)
\end{minipage} & \begin{minipage}[b]{\linewidth}\raggedright
GMC
\end{minipage} & \begin{minipage}[b]{\linewidth}\centering
\(\supset\)
\end{minipage} & \begin{minipage}[b]{\linewidth}\raggedright
clump / filament
\end{minipage} & \begin{minipage}[b]{\linewidth}\centering
\(\supset\)
\end{minipage} & \begin{minipage}[b]{\linewidth}\raggedright
dense core
\end{minipage} & \begin{minipage}[b]{\linewidth}\centering
\(\to\)
\end{minipage} & \begin{minipage}[b]{\linewidth}\raggedright
protostar + disk
\end{minipage} \\
\midrule\noalign{}
\endhead
\bottomrule\noalign{}
\endlastfoot
size & \(10\)--\(100\ \mathrm{pc}\) & & \(20\)--\(100\ \mathrm{pc}\) & &
clump: \(0.3\)--\(3\ \mathrm{pc}\) filament: \(1\)--\(10\ \mathrm{pc}\)
\(\times\) \(0.1\ \mathrm{pc}\) & & radius
\(0.03\)--\(0.1\ \mathrm{pc}\) diameter \(0.06\)--\(0.2\ \mathrm{pc}\) &
& envelope: \(10^3\)--\(10^4\ \mathrm{AU}\) disk:
\(30\)--\(300\ \mathrm{AU}\) \\
mass & \(10^3\)--\(10^6\,M_\odot\) & & \(10^4\)--\(10^6\,M_\odot\) & &
\(10\)--\(10^3\,M_\odot\) & & \(0.5\)--\(10\,M_\odot\) & & protostar:
\(0.1\)--a few \(M_\odot\) disk: \(10^{-3}\)--\(10^{-1}\,M_\odot\) \\
number density & \(1\)--\(10^2\ \mathrm{cm^{-3}}\) & &
\(10^2\)--\(10^3\ \mathrm{cm^{-3}}\) & &
\(10^3\)--\(10^5\ \mathrm{cm^{-3}}\) & &
\(10^4\)--\(10^6\ \mathrm{cm^{-3}}\) & & envelope:
\(10^5\)--\(10^8\ \mathrm{cm^{-3}}\) disk midplane can be much higher \\
temperature & \(10^2\)--\(10^4\ \mathrm{K}\) & &
\(10\)--\(20\ \mathrm{K}\) & & \(10\)--\(20\ \mathrm{K}\) & &
\(10\ \mathrm{K}\) & & envelope: \(10\)--\(100\ \mathrm{K}\) inner disk
can be much hotter \\
physics & large-scale ISM environment; only its coldest, densest parts
are relevant for star formation & & cold, self-gravitating molecular
reservoir & & fragmentation into local star-forming sites; often
produces groups or clusters & & smallest prestellar object plausibly
associated with one star system or small multiple & & gravity,
accretion, angular momentum transport, jets/outflows, and eventual
planet formation \\
\end{longtable}

Table 2: The structures relevant to the first few stages of star
formation. The symbols should be read horizontally as ``contains'',
``includes'', ``contains'', and ``evolves into''. Thus an ISM complex
contains GMCs, a GMC contains clumps and filaments, and these contain
dense cores. All values are only order-of-magnitude estimates; real
clouds are irregular, turbulent, magnetized, and far from spherical.

Additional notes - The \textbf{dense core} is the first entry here that
is often associated with the formation of a single star system (or a
small multiple system). - The last step is written with an arrow because
a collapsing dense core evolves into a \textbf{protostar} plus disk. The
protostar is dense shrouded core that is not yet visible from the
outside (it's covered in dense material and not luminous enough to shine
through).

Here is a recap (with some new tools sprinkled in) of the physics
relevant to the collapse.

\paragraph{Astrophysics of collapse
onset}\label{astrophysics-of-collapse-onset}

A GMC is self-gravitating but not in simple free fall. It is cold,
molecular, turbulent, and magnetized, so gravity is opposed not only by
thermal pressure but also by turbulent motions and magnetic stresses. A
useful diagnostic is the virial parameter, \[
\alpha_{\rm vir}=\frac{5\sigma^2 R}{GM},
\] with \(\alpha_{\rm vir}\sim 1\) indicating that self-gravity is
dynamically important. In modern star-formation theory, turbulence plays
a double role: on large scales it helps prevent the whole cloud from
collapsing monolithically, while on smaller scales it creates overdense
clumps and filaments that can become locally unstable.

Once a dense region becomes unstable, the natural collapse timescale is
the \textbf{free-fall time}, \[
t_{\rm ff}=\sqrt{\frac{3\pi}{32G\rho}}.
\] Because \(t_{\rm ff}\propto 1/\sqrt{\rho}\), denser substructures
collapse faster than their parent cloud, which is one reason star
formation proceeds hierarchically. \textbf{Fragmentation} occurs because
the collapsing gas is not uniform: turbulence seeds density structure,
and regions whose mass exceeds the local Jeans mass can collapse
separately. While the gas remains approximately isothermal,
fragmentation is comparatively easy; later heating raises the effective
pressure and makes further small-scale fragmentation harder.

After a central object forms, star formation becomes largely an
\textbf{accretion problem}. Gas does not fall directly onto the
protostar from all directions, because even a slowly rotating core
contains too much angular momentum for purely radial infall. (Recall
that the angular momentum \(L\propto MR^2 \omega\) must be conserved, so
if an initially large \(R\) decreases by a factor of 100 (stage
\(1\to2\) or \(3\to4\)) then \(\omega\) must increase by a factor of
\((100)^2=10,000\).) Instead, the infalling gas settles into a disk, and
the protostar grows as mass accretes through that disk. A useful
estimate for the luminosity in this phase is the accretion luminosity,
\[
L_{\rm acc}\sim \frac{GM\dot M}{R}
\] which explains why young protostars can be luminous before sustained
hydrogen fusion begins. (Here \(\dot M=dM/dt\) is the time rate of
change of the mass of the protostar.)

The central theoretical difficulty is angular momentum transport. From
dense core to young star, the angular momentum must decrease by several
orders of magnitude. Part of that transport happens already in the
infalling envelope through primarily magnetized disk winds/jets, which
remove angular momentum vertically and allow matter to spiral inward. In
many modern treatments, magnetic torques and magnetized outflows are not
side effects but one of the main reasons accretion can continue at all.

The temperature rise in the protostellar core and later the pre-main
sequence (PMS) star/object is initially not caused by fusion, but by
compressional heating. Early in collapse, the core is optically thin
enough that dust thermal emission can efficiently radiate away the
released gravitational energy, so the collapse is nearly isothermal. As
density rises, the gas and dust become optically thick, radiative
cooling can no longer keep pace with \(PdV\) heating (do you know what
this means??), and the central temperature begins to increase. This
produces the first \textbf{hydrostatic} core:
\href{https://en.wikipedia.org/wiki/Hydrostatic_equilibrium}{\textbf{Hydrostatic
equilibrium}} is defined as when a liquid or other ``plastic''
(deformable) body is static (not changing shape) due to a balance
between forces with one of the forces arising from pressure. In this
case, the force of gravity pulling inward is balanced by a pressure
``gradient'' with the pressure (force per unit area) increasing as you
get deeper into the center.

When the central temperature reaches roughly \(2000\) \texttt{K},
molecular hydrogen begins to dissociate; that dissociation absorbs
energy, softens the equation of state (read: makes the gas/liquid easier
to compress), and triggers a second collapse toward the true
protostellar core.

    \subsection{\texorpdfstring{The
\href{https://en.wikipedia.org/wiki/Kelvin–Helmholtz_mechanism}{Kelvin-Helmholtz}
phase}{The Kelvin-Helmholtz phase}}\label{the-kelvin-helmholtz-phase}

During phase 3, the density of the material in the protostar increases
to the point that it becomes opaque and radiation becomes trapped. This
causes the core inner region to heat up from 100 \texttt{K} to 10,000
\texttt{K}. Meanwhile, the outer regions of the dense part of the cloud
tens of au further out are still thin and transparent so the outer
regions are still around the 100 \texttt{K} scale. A photosphere becomes
sharp as the ``edge'' of the dense region becomes more pronounced and a
\textbf{pre-main sequence} star (PMS star) has formed. The luminosity of
the PMS star is a power (= energy per unit time) and that power is
supplied by the gravitational energy released during contraction. This
is called
\href{https://en.wikipedia.org/wiki/Kelvin–Helmholtz_mechanism}{\textbf{Kelvin-Helmholtz}}
contraction.

Gravitation will continue to be the primary source of luminosity until
the core reaches the 10 \texttt{MK} \(= 10^7\) \texttt{K} needed to
initiate proton-proton fusion, which is achieved only in stage 6. Thus
during stages 3 through 5, the protostar is in the KH phase. Throughout
this phase, the internal pressure is not yet enough to balance the
gravitational attraction, so the star continues to shrink.

 
\subsubsection{Problem 8}\label{problem-8}

We can see from the HR diagram that during stage 3, the radius shrinks
only moderately (the track is roughly parallel to the \(100\ R_\odot\)
line) and the luminosity increases only moderately (the track is roughly
horizontal) as well. Let's use these two observations to estimate the
time this phase takes by assuming the luminosity \(L\) and radius \(R\)
stay roughly constant. Since luminosity is the rate of change of energy,
we need the amount of energy available. By definition, in the
Kelvin-Helmholtz phase, this energy is coming from gravitational
attraction, so we will estimate it from the gravitational potential
energy.

\begin{enumerate}
\def\labelenumi{\alph{enumi}.}
\tightlist
\item
  Use dimensional analysis to show that the gravitational potential
  energy of a star of radius \(R\) and mass \(M\) that formed from
  particles that were initially very widely separated should be of
  order\\
  \[
  U_g \sim \frac{GM^2}{R}\tag{7}
  \]
\end{enumerate}

The intuition behind this is that all the particles that make up this
star were separated by parsec distances, so the gravitational energy
between them was very small. As they come in to a distance of radius
\(R\) to make up a star, we get a lot of energy out. This energy is what
heats up the star. If we do this calculation carefully, we will need
calculus, but the only thing that changes is the numerical coefficient
in the relation you derived from dimensional analysis. (The numerical
coefficient is \(\frac35\).)

\begin{enumerate}
\def\labelenumi{\alph{enumi}.}
\setcounter{enumi}{1}
\item
  Calculate the gravitational potential energy for a solar mass star
  from your estimate.
\item
  The luminosity of the sun \(L_\odot=3.8\times 10^{26}\) \texttt{W}.
  Since this is a rate of energy consumption (\emph{i.e.} power), we can
  estimate the time it takes to consume any amount of energy: time
  \textasciitilde{} energy/power. Use this and the HR diagram above to
  estimate the order of magnitude of the amount of time a solar mass
  protostar will be in the Kelvin-Helmholtz contraction phase. (Compare
  this to table 1.)
\end{enumerate}
 

     
\subsubsection{Answer 8}\label{answer-8}

\begin{enumerate}
\def\labelenumi{\alph{enumi}.}
\item
  From Newton's Second Law \(F=ma\), the dimensions of force are
  \([F]=[m][a]=M \cdot L \cdot T^{-2}\). From Newton's Law of Universal
  Gravitation \(F=Gm_1m_2/r^2\) this is also
  \([G]\cdot M^2\cdot L^{-2}\), so the gravitational constant has
  dimensions \([G]=M^{-1}\cdot L^3 \cdot T^{-2}\). From the kinetic
  energy \(K=\frac12 mv^2\) (or Einstein's relation \(E=mc^2\)), energy
  has dimensions \([U]=M\cdot L^2 \cdot T^{-2}\). To match the
  mass-dependance, we need \(U\propto GM^2\). This combination has
  remaining dimensions \(M\cdot L^3\cdot T^{-2}\), so to cancel one
  power of the length, we need \(U\propto GM^2/R\), as claimed.
\item
  See below.
\item
  At point 4, the luminosity \(L\approx 200\ L_\odot\). Below, we plug
  this into the estimate for the time scale \(t \sim U/L\).
\end{enumerate}
 

    \begin{tcolorbox}[breakable, size=fbox, boxrule=1pt, pad at break*=1mm,colback=cellbackground, colframe=cellborder]
\prompt{In}{incolor}{7}{\boxspacing}
\begin{Verbatim}[commandchars=\\\{\}]
\PY{k+kn}{from}\PY{+w}{ }\PY{n+nn}{astropy}\PY{n+nn}{.}\PY{n+nn}{constants}\PY{+w}{ }\PY{k+kn}{import} \PY{n}{G}\PY{p}{,} \PY{n}{R\PYZus{}sun}\PY{p}{,} \PY{n}{L\PYZus{}sun}

\PY{n+nb}{print}\PY{p}{(}\PY{l+s+sa}{f}\PY{l+s+s2}{\PYZdq{}}\PY{l+s+se}{\PYZbs{}n}\PY{l+s+s2}{|   ANSWER 8}\PY{l+s+se}{\PYZbs{}n}\PY{l+s+s2}{|}\PY{l+s+s2}{\PYZdq{}}\PY{p}{)}

\PY{n}{\PYZus{}U} \PY{o}{=} \PY{n}{G}\PY{o}{*}\PY{n}{M\PYZus{}sun}\PY{o}{*}\PY{o}{*}\PY{l+m+mi}{2}\PY{o}{/}\PY{n}{R\PYZus{}sun}

\PY{n}{\PYZus{}L}\PY{o}{=} \PY{l+m+mi}{200}\PY{o}{*}\PY{n}{L\PYZus{}sun}

\PY{n}{\PYZus{}t} \PY{o}{=} \PY{n}{\PYZus{}U}\PY{o}{/}\PY{n}{\PYZus{}L}

\PY{n+nb}{print}\PY{p}{(}\PY{l+s+sa}{f}\PY{l+s+s2}{\PYZdq{}}\PY{l+s+s2}{|   b. With a mass of }\PY{l+s+si}{\PYZob{}}\PY{n}{M\PYZus{}sun}\PY{l+s+si}{:}\PY{l+s+s2}{.1e}\PY{l+s+si}{\PYZcb{}}\PY{l+s+s2}{, and a radius of }\PY{l+s+si}{\PYZob{}}\PY{n}{R\PYZus{}sun}\PY{l+s+si}{:}\PY{l+s+s2}{.1e}\PY{l+s+si}{\PYZcb{}}\PY{l+s+s2}{, we find an energy of}\PY{l+s+se}{\PYZbs{}n}\PY{l+s+s2}{|      }\PY{l+s+si}{\PYZob{}}\PY{n}{\PYZus{}U}\PY{o}{.}\PY{n}{to}\PY{p}{(}\PY{n}{u}\PY{o}{.}\PY{n}{J}\PY{p}{)}\PY{l+s+si}{:}\PY{l+s+s2}{.1e}\PY{l+s+si}{\PYZcb{}}\PY{l+s+s2}{.}\PY{l+s+se}{\PYZbs{}n}\PY{l+s+s2}{|}\PY{l+s+s2}{\PYZdq{}}\PY{p}{)}

\PY{n+nb}{print}\PY{p}{(}\PY{l+s+sa}{f}\PY{l+s+s2}{\PYZdq{}}\PY{l+s+s2}{|   c. At a luminosity of }\PY{l+s+si}{\PYZob{}}\PY{n}{\PYZus{}L}\PY{l+s+si}{:}\PY{l+s+s2}{.1e}\PY{l+s+si}{\PYZcb{}}\PY{l+s+s2}{ (i.e. Joules per second), equivalent to that of}\PY{l+s+se}{\PYZbs{}n}\PY{l+s+s2}{|      200 Suns, this translates to a time\PYZhy{}scale of }\PY{l+s+si}{\PYZob{}}\PY{n}{\PYZus{}t}\PY{o}{.}\PY{n}{decompose}\PY{p}{(}\PY{p}{)}\PY{l+s+si}{:}\PY{l+s+s2}{.1e}\PY{l+s+si}{\PYZcb{}}\PY{l+s+s2}{. In more intuitive}\PY{l+s+se}{\PYZbs{}n}\PY{l+s+s2}{|      units, this is }\PY{l+s+si}{\PYZob{}}\PY{n}{\PYZus{}t}\PY{o}{.}\PY{n}{to}\PY{p}{(}\PY{n}{u}\PY{o}{.}\PY{n}{yr}\PY{p}{)}\PY{l+s+si}{:}\PY{l+s+s2}{.1e}\PY{l+s+si}{\PYZcb{}}\PY{l+s+s2}{. According to tabel 1, this is exactly the right}\PY{l+s+se}{\PYZbs{}n}\PY{l+s+s2}{|      dex for phase 3.}\PY{l+s+se}{\PYZbs{}n}\PY{l+s+s2}{\PYZdq{}}\PY{p}{)}
\end{Verbatim}
\end{tcolorbox}

    \begin{Verbatim}[commandchars=\\\{\}]

|   ANSWER 8
|
|   b. With a mass of 2.0e+30 kg, and a radius of 7.0e+08 m, we find an energy
of
|      3.8e+41 J.
|
|   c. At a luminosity of 7.7e+28 W (i.e. Joules per second), equivalent to that
of
|      200 Suns, this translates to a time-scale of 5.0e+12 s. In more intuitive
|      units, this is 1.6e+05 yr. According to tabel 1, this is exactly the
right
|      dex for phase 3.

    \end{Verbatim}

    From this calculation, we confirm that we understand the physics
responsible for the contraction and luminosity of the Kelvin-Helmholtz
phase.

\subsubsection{Stages 4 and 5: T Tauri
phase}\label{stages-4-and-5-t-tauri-phase}

Checking against our table, we confirm that the stage 3-to-4 handover
point has \(T=3000\) \texttt{K} and radius \(100\ R_\odot\), so the
protostar is \textasciitilde100 times the size of the main sequence star
it will finally settle down to. At this stage, we start down the
\href{https://en.wikipedia.org/wiki/Hayashi_track}{\textbf{Hayashi
track}}. The Hayashi limit is defined as the line across which a star
goes from being fully convective as opposed to forming a radiation zone.
To the left of the line, internal temperatures are so high that
photoionization becomes significant and energy is transported primarily
through radiation. To the right of the line, convection is so great that
large quantities of thermal energy are rapidly transported to the
surfact. Hayashi showed, that being even a small amount to the right of
this line, the cooling effect on the interior is so great, that the star
quickly adujsts back to the \textbf{Hayashi temperature} (a.k.a. the
Hayashi limit): The lowest temperature that a fully convective star will
be hydrodynamically stable.

Thus, in this phase of protostellar evolution, the star is fully
convective (a radiation zone has not yet formed around the core) and is
continuing to contract under the influence of gravity. It's surface
temperature stays roughly the same, so it's luminosity drops (why?
\ldots{} at this point, you should be instinctively checking statements
I make against things you already understand), and the track is amost
vertical.

The name of this pre-main sequence (PMS) object is a
\href{https://en.wikipedia.org/wiki/T_Tauri_star}{\textbf{T Tauri star}}
in honor of the first PMS object observed in the constellation Taurus
(shown at left). It is a \textbf{variable} star meaning that its
luminosity is changing in time. Thick convection zones transport energy
to the surface which often results in violent eruptions. During this
phase, the protostar gives rise to the strong protostellar winds
responsible for clearing the protostellar system of interplanetary gas.
Early in this phase, the winds are collimated into
\href{https://en.wikipedia.org/wiki/Astrophysical_jet}{\textbf{jets}}
resulting in
\href{https://en.wikipedia.org/wiki/Bipolar_outflow}{\textbf{bipolar
outflow}} from the poles and perpendicular to the protoplanetary disk.
Often, these streams of ejected material collide with the protostellar
medium around it making large and visible shock fronts called
\href{https://en.wikipedia.org/wiki/Herbig–Haro_object}{\textbf{Herbig-Haro
objects}}. Herbig-Haro object HH30 is shown in the \texttt{gif} to the
right.

Here is a link to a
\href{https://www.youtube.com/watch?v=8nCW7419FyU&t=63s}{recent
simulation} by \href{https://starforge.space/index.html}{Starforge}
demonstrating star formation from an interstellar cloud. (See also
\href{https://www.youtube.com/watch?v=PUyE3j0aoMw}{this}.) In the
simulation we see many instances of mass outflows of the T Tauri stars.

 
\subsubsection{Problem 9.}\label{problem-9.}

Watch the Starforge simulation linked above. When do the first
protostars enter the T Tauri phase? At some point near the end, there is
a large flash. When does this flash happen and what do you think is
happening then?
 

     
\subsubsection{Answer 9.}\label{answer-9.}

The first T Tauri-phase stars occur around 01:00 and the supernova is
around 03:00.
 

    \subsection{Stage 6}\label{stage-6}

As the T Tauri star progresses down the Hayashi track with a
roughly-constant surface temperature the luminosity is decreasing in
accordance with the Stefan-Boltzmann law for a shrinking object. The
heat generated by this contraction acts to raise the core temperature.
This increase, in turn, causes the opacity of the core to drop below a
point where radiative energy transportation becomes competitive: a
radiation zone starts to form and then grows in size. This transfers us
off the Hayashi track onto the
\href{https://en.wikipedia.org/wiki/Henyey_track}{\textbf{Henyey
track}}.

At this point, we have a radiative zone inside of a convective outer
region. Although the Kelvin-Helmholtz contraction is still ongoing, core
temperature and pressure have rising to the point that hydrogen fusion
begins at 10 \texttt{MK} \(=\) 10,000,000 \texttt{K}. As contraction
proceeds the luminosity begins to transition between
gravitationally-powered to fusion-powered.

Although the protostar has become a star, the internal pressure is still
not balancing the gravitational attraction, and contraction continues
until stage 7, where the star finally reaches
\href{https://en.wikipedia.org/wiki/Hydrostatic_equilibrium}{\textbf{hydrostatic
equilibrium}}.

\subsection{Stage 7}\label{stage-7}

At the end of stage 6, the protostar has become a star, but the internal
pressure is still not balancing the gravitational attraction, and
contraction continues until stage 7, where the star finally reaches
\href{https://en.wikipedia.org/wiki/Hydrostatic_equilibrium}{\textbf{hydrostatic
equilibrium}}. This is the final condition needed to define a
\textbf{zero-age main sequence} star.

What defines a \textbf{star on the main sequence} is that * it ``burns''
hydrogen specifically (\emph{i.e.} fuses it into helium), and\\
* it is in hydrostatic equilibrium.

Eventually all such stars must burn through their hydrogen. The rate at
which they do so is proportional to their luminosity (see problem 8).

 
\subsubsection{Problem 10}\label{problem-10}

Estimate the time a solar-mass star with spectral type G will remain on
the main sequence. The fusion of protons in the sun generates
\(\Delta E_{pp}=1\times 10^{-12}\) \texttt{J} per proton.

\begin{enumerate}
\def\labelenumi{\alph{enumi}.}
\item
  Calculate the number of protons in the Sun by assuming 75\% of it is
  made of protons.
\item
  Assume 20\% of this is available for fusion at a rate sustaining the
  solar luminosity \(L_\odot\) of the Sun, and use this to estimate the
  time this star will remain on the main sequence.
\end{enumerate}
 

    \begin{tcolorbox}[breakable, size=fbox, boxrule=1pt, pad at break*=1mm,colback=cellbackground, colframe=cellborder]
\prompt{In}{incolor}{8}{\boxspacing}
\begin{Verbatim}[commandchars=\\\{\}]
\PY{n+nb}{print}\PY{p}{(}\PY{l+s+sa}{f}\PY{l+s+s2}{\PYZdq{}}\PY{l+s+se}{\PYZbs{}n}\PY{l+s+s2}{|   ANSWER 10}\PY{l+s+se}{\PYZbs{}n}\PY{l+s+s2}{|}\PY{l+s+s2}{\PYZdq{}}\PY{p}{)}

\PY{n}{\PYZus{}M} \PY{o}{=} \PY{l+m+mi}{3}\PY{o}{/}\PY{l+m+mi}{4}\PY{o}{*}\PY{n}{M\PYZus{}sun}

\PY{n}{\PYZus{}N} \PY{o}{=} \PY{n}{\PYZus{}M}\PY{o}{/}\PY{n}{m\PYZus{}p}

\PY{n+nb}{print}\PY{p}{(}\PY{l+s+sa}{f}\PY{l+s+s2}{\PYZdq{}}\PY{l+s+s2}{|   a. 3/4 of the mass of the Sun is about }\PY{l+s+si}{\PYZob{}}\PY{n}{\PYZus{}M}\PY{l+s+si}{:}\PY{l+s+s2}{.1e}\PY{l+s+si}{\PYZcb{}}\PY{l+s+s2}{ giving }\PY{l+s+si}{\PYZob{}}\PY{n}{\PYZus{}N}\PY{l+s+si}{:}\PY{l+s+s2}{.1e}\PY{l+s+si}{\PYZcb{}}\PY{l+s+s2}{ protons.}\PY{l+s+se}{\PYZbs{}n}\PY{l+s+s2}{|}\PY{l+s+s2}{\PYZdq{}}\PY{p}{)}

\PY{n}{\PYZus{}Neff} \PY{o}{=} \PY{l+m+mf}{.2}\PY{o}{*}\PY{n}{\PYZus{}N}

\PY{n}{\PYZus{}E} \PY{o}{=} \PY{n}{\PYZus{}N}\PY{o}{*}\PY{l+m+mf}{1e\PYZhy{}12}\PY{o}{*}\PY{n}{u}\PY{o}{.}\PY{n}{J}

\PY{n}{\PYZus{}t} \PY{o}{=} \PY{n}{\PYZus{}E}\PY{o}{/}\PY{n}{L\PYZus{}sun}

\PY{n+nb}{print}\PY{p}{(}\PY{l+s+sa}{f}\PY{l+s+s2}{\PYZdq{}}\PY{l+s+s2}{|   b. 20\PYZpc{} of these protons is about }\PY{l+s+si}{\PYZob{}}\PY{n}{\PYZus{}Neff}\PY{l+s+si}{:}\PY{l+s+s2}{.1e}\PY{l+s+si}{\PYZcb{}}\PY{l+s+s2}{ giving and energy of }\PY{l+s+si}{\PYZob{}}\PY{n}{\PYZus{}E}\PY{l+s+si}{:}\PY{l+s+s2}{.1e}\PY{l+s+si}{\PYZcb{}}\PY{l+s+s2}{.}\PY{l+s+se}{\PYZbs{}n}\PY{l+s+s2}{|      At a rate of }\PY{l+s+si}{\PYZob{}}\PY{n}{L\PYZus{}sun}\PY{l+s+si}{:}\PY{l+s+s2}{.1e}\PY{l+s+si}{\PYZcb{}}\PY{l+s+s2}{, this is a lifetime }\PY{l+s+si}{\PYZob{}}\PY{n}{\PYZus{}t}\PY{l+s+si}{:}\PY{l+s+s2}{.1e}\PY{l+s+si}{\PYZcb{}}\PY{l+s+s2}{ or }\PY{l+s+si}{\PYZob{}}\PY{n}{\PYZus{}t}\PY{o}{.}\PY{n}{to}\PY{p}{(}\PY{n}{u}\PY{o}{.}\PY{n}{Gyr}\PY{p}{)}\PY{l+s+si}{:}\PY{l+s+s2}{.1e}\PY{l+s+si}{\PYZcb{}}\PY{l+s+s2}{. }\PY{l+s+se}{\PYZbs{}n}\PY{l+s+s2}{|}\PY{l+s+s2}{\PYZdq{}}\PY{p}{)}
\end{Verbatim}
\end{tcolorbox}

    \begin{Verbatim}[commandchars=\\\{\}]

|   ANSWER 10
|
|   a. 3/4 of the mass of the Sun is about 1.5e+30 kg giving 8.9e+56 protons.
|
|   b. 20\% of these protons is about 1.8e+56 giving and energy of 8.9e+44 J.
|      At a rate of 3.8e+26 W, this is a lifetime 2.3e+18 J / W or 7.4e+01 Gyr.
|
    \end{Verbatim}

    \subsubsection{Mass-Luminosity/Radius
relations}\label{mass-luminosityradius-relations}

\subsubsection{Problem 11}\label{problem-11}

Consider the diagram at right (which I have borrowed from Carroll and
Ostlie).

\begin{enumerate}
\def\labelenumi{\alph{enumi}.}
\item
  Describe what you see.
\item
  Use it to construct a formula for \(L\) as a function of \(M\).
\end{enumerate}

\subsubsection{Answer 11}\label{answer-11}

\begin{enumerate}
\def\labelenumi{\alph{enumi}.}
\tightlist
\item
  It appears to be a diagram of the luminosity of a collection of stars
  plotted against the mass of those stars.
\end{enumerate}

\begin{itemize}
\tightlist
\item
  From the legend, I gather that these stars are in various types of
  binary systems, which is presumably how they determined the masses.
\item
  I also notice that all the stars are very different in mass ranging
  from \(0.1\ M_\odot\) to \(10^{1.25}\ M_\odot \approx 18\ M_\odot\).
  This is confirmed by ``OB systems'' label in the legend as OB stars
  are very hot an luminous. Even across such large mass ranges, the
  relation is quite linear.
\end{itemize}

\begin{enumerate}
\def\labelenumi{\alph{enumi}.}
\setcounter{enumi}{1}
\tightlist
\item
  Taking the run from \(-0.5\) to \(1.25\) and the rise from \(-2\) to
  \(4.5\), I'm getting a slope of \(\alpha=6.5/1.75\approx 3.7\). The
  relation is, therefore \(L/L_\odot = \left( M/M_\odot\right)^\alpha\)
  with this value for \(\alpha\).
\end{enumerate}


    This is called the
\href{https://en.wikipedia.org/wiki/Mass\%E2\%80\%93luminosity_relation}{\textbf{mass-luminosity
relation}}. Actually, there are multiple such relations that are all of
the form \[
L \propto M^\alpha \tag{8}
\] for some constant \(\alpha\). The value of \(\alpha\) quoted in the
book is \(\alpha\approx 4\). I usually see this quoted to be
\(\alpha\approx 3.5\). Later in the course when we look at the
particulars of stellar evolution, we may come back and tweak this later.
I will stick with \(\alpha \approx 3.7\), but let me briefly sketch the
reason such a relation is not entirely unexpected.

We will use three main ideas that may or may not be good
approximations/models: * Hydrostatic equilibrium: pressure gradients are
what are holding up the star against gravitational collabpse. This is
reasonable because it is one of the two defining conditions for a
main-sequence star. * Ideal gas equation of state. This is a bit iffy,
but if we can think of a star as a giant ball of ideal gas, then
\(PV=Nk_BT\). * Radiative diffusion: The energy transport out of the
core is carried by radiation. This is very iffy for very massive or very
low-mass stars, as we will discuss later.

The first assumption lets us say that the pressure
\(P=\textrm{energy}/V\) at the core is balancing \(U_g=GM^2/R\), thus \[
P\sim \frac{GM^2}{R^4}
\] From the second assumption, this is equal to
\(P=\frac{\rho}{m_p} k_BT\) where \(\rho\) is the mass density, so
\(\frac{\rho}{m_p}\) is the number density. This gives the core
temperature as \[
k_B T_c\sim \frac{Gm_p M}{R}
\] The final assumption allow us to say that the energy transport out of
the core has a flux given by the Stefan-Boltzmann law
\(F=\sigma_B T^4\). Unfortunately, we cannot blindly plug in SB in the
form \(L\sim R^2 T^4\) because this formula was appropriate for the
surface of the star and using it would only be correct if
\(T_{eff}=T_c\) which is very much wrong, as we can see from table 1. To
proceed this way, we need to talk about \emph{opacity} which is how
not-transparent the material is to radiation. What lowers the
temperature as we go outward is precisly the opacity \(\kappa\) that we
are missing. It is usually a reasonable assumption to take the opacity
\(\kappa\) to be constant, but this factor has dimensions
\([\kappa]=L^2/M\), so its existence will change the dimensional
analysis, because the combination \(\frac{\kappa M}{R^2}=\kappa \rho R\)
is dimensionless. Specifically, the flux is modified by opacity to \[
F\sim \frac{\sigma_B T_c^4}{\kappa \rho R}
\] where we are taking \(\kappa\) and \(\rho\) to be constant. Note that
while this is a reasonable guess (\emph{e.g.} as opacity increases, the
flux decreases), a derivation of this result is needed in order to
justify this formula. (It is not a difficult derivation, but this packet
is already too long.)

In any case, if we accept this modification, we can proceed as before,
as the luminosity now follows \[
L \sim R^2F\sim \frac{\sigma_B R T_c^4}{\kappa \rho} \sim \frac{\sigma_B R^4}{\kappa M}\left(\frac{Gm_p M}{R}\right)^4 \propto \frac{M^3}{\kappa}
\] Notice that it is very convenient that the radius has canceled since
we do not (yet) know the relation (if any) between mass and radius.

As we already discussed in Problem 10, the luminosity represents the
amount of energy emitted per unit time (the rate of burning), so the
lifetime \(t\) of a star is inversely proportional to the luminosity
\(t\propto 1/L\). Since there is more fuel to burn when there is more
mass, it is also proportional to the mass \(t\propto M\). Taken together
with the mass-luminosity relation (8), we get that the lifetime of a
star is inversely proportional to the cube of the mass of the star: \[
                t\propto \frac1{M^{\alpha -1}}  \tag{8}
\] This means that the more massive the star, the shorter it stays on
the main sequence branch.

 
\subsubsection{Problem 12}\label{problem-12}

An oft-cited estimate for the total lifetime of the sun is
\(t=1\times 10^{10}\) \texttt{yr} \(10\) \texttt{Gyr}. Use this to
estimate the lifetime of a 15-solar-mass main sequence star. What about
a \(60M_\odot\) star?
 

    \begin{tcolorbox}[breakable, size=fbox, boxrule=1pt, pad at break*=1mm,colback=cellbackground, colframe=cellborder]
\prompt{In}{incolor}{9}{\boxspacing}
\begin{Verbatim}[commandchars=\\\{\}]
\PY{n+nb}{print}\PY{p}{(}\PY{l+s+sa}{f}\PY{l+s+s2}{\PYZdq{}}\PY{l+s+se}{\PYZbs{}n}\PY{l+s+s2}{|   ANSWER 12}\PY{l+s+se}{\PYZbs{}n}\PY{l+s+s2}{|}\PY{l+s+s2}{\PYZdq{}}\PY{p}{)}
\PY{n}{t\PYZus{}sun} \PY{o}{=} \PY{l+m+mi}{10\PYZus{}000\PYZus{}000\PYZus{}000}
\PY{n}{t\PYZus{}15} \PY{o}{=} \PY{l+m+mi}{15}\PY{o}{*}\PY{o}{*}\PY{p}{(}\PY{o}{\PYZhy{}}\PY{l+m+mi}{3}\PY{p}{)}
\PY{n}{t\PYZus{}60} \PY{o}{=} \PY{l+m+mi}{60}\PY{o}{*}\PY{o}{*}\PY{p}{(}\PY{o}{\PYZhy{}}\PY{l+m+mi}{3}\PY{p}{)}
\PY{n+nb}{print}\PY{p}{(}\PY{l+s+sa}{f}\PY{l+s+s2}{\PYZdq{}}\PY{l+s+s2}{|   A 15\PYZhy{}solar mass star has a lifetime that is }\PY{l+s+si}{\PYZob{}}\PY{n+nb}{round}\PY{p}{(}\PY{l+m+mi}{1}\PY{o}{/}\PY{n}{t\PYZus{}15}\PY{p}{)}\PY{l+s+si}{\PYZcb{}}\PY{l+s+s2}{ times shorter than our Sun, }\PY{l+s+se}{\PYZbs{}n}\PY{l+s+s2}{|   or }\PY{l+s+si}{\PYZob{}}\PY{n+nb}{int}\PY{p}{(}\PY{n+nb}{round}\PY{p}{(}\PY{n}{t\PYZus{}15}\PY{o}{*}\PY{n}{t\PYZus{}sun}\PY{o}{/}\PY{l+m+mf}{1e6}\PY{p}{)}\PY{p}{)}\PY{l+s+si}{\PYZcb{}}\PY{l+s+s2}{ million years.}\PY{l+s+se}{\PYZbs{}n}\PY{l+s+s2}{|}\PY{l+s+s2}{\PYZdq{}}\PY{p}{)}
\PY{n+nb}{print}\PY{p}{(}\PY{l+s+sa}{f}\PY{l+s+s2}{\PYZdq{}}\PY{l+s+s2}{|   A 60\PYZhy{}solar mass star has a lifetime that is }\PY{l+s+si}{\PYZob{}}\PY{n+nb}{round}\PY{p}{(}\PY{l+m+mi}{1}\PY{o}{/}\PY{n}{t\PYZus{}60}\PY{p}{)}\PY{l+s+si}{\PYZcb{}}\PY{l+s+s2}{ times shorter than our Sun, }\PY{l+s+se}{\PYZbs{}n}\PY{l+s+s2}{|   or }\PY{l+s+si}{\PYZob{}}\PY{n+nb}{int}\PY{p}{(}\PY{n+nb}{round}\PY{p}{(}\PY{n}{t\PYZus{}60}\PY{o}{*}\PY{n}{t\PYZus{}sun}\PY{o}{/}\PY{l+m+mf}{1e3}\PY{p}{,}\PY{o}{\PYZhy{}}\PY{l+m+mi}{1}\PY{p}{)}\PY{p}{)}\PY{l+s+si}{\PYZcb{}}\PY{l+s+s2}{ thousand years.}\PY{l+s+se}{\PYZbs{}n}\PY{l+s+s2}{\PYZdq{}}\PY{p}{)}
\end{Verbatim}
\end{tcolorbox}

    \begin{Verbatim}[commandchars=\\\{\}]

|   ANSWER 12
|
|   A 15-solar mass star has a lifetime that is 3375 times shorter than our Sun,
|   or 3 million years.
|
|   A 60-solar mass star has a lifetime that is 216000 times shorter than our
Sun,
|   or 50 thousand years.

    \end{Verbatim}

     
\subsubsection{Problem 13}\label{problem-13}

Go back to table 1 and add up all the time it takes to form a star, that
is the time from the beginning of stage 1 until the end of stage 6.
Compare this to the time the star stays on the main sequence. What
percentage of the latter time is the former?
 

    \begin{tcolorbox}[breakable, size=fbox, boxrule=1pt, pad at break*=1mm,colback=cellbackground, colframe=cellborder]
\prompt{In}{incolor}{10}{\boxspacing}
\begin{Verbatim}[commandchars=\\\{\}]
\PY{n+nb}{print}\PY{p}{(}\PY{l+s+sa}{f}\PY{l+s+s2}{\PYZdq{}}\PY{l+s+se}{\PYZbs{}n}\PY{l+s+s2}{|   ANSWER 13}\PY{l+s+se}{\PYZbs{}n}\PY{l+s+s2}{|}\PY{l+s+s2}{\PYZdq{}}\PY{p}{)}
\PY{n+nb}{print}\PY{p}{(}\PY{l+s+sa}{f}\PY{l+s+s2}{\PYZdq{}}\PY{l+s+s2}{|   After the main sequence at 10 Gyr, stages 5 and 6 are the longest at 10 Myr;}\PY{l+s+se}{\PYZbs{}n}\PY{l+s+s2}{|   each is smaller by a factor of 1000. As a fraction, this means that a star is on}\PY{l+s+se}{\PYZbs{}n}\PY{l+s+s2}{|   the main sequence 1/(1+2/1000) of the time, or percentage\PYZhy{}wise}\PY{l+s+se}{\PYZbs{}n}\PY{l+s+s2}{|   }\PY{l+s+si}{\PYZob{}}\PY{l+m+mi}{100}\PY{o}{/}\PY{p}{(}\PY{l+m+mi}{1}\PY{o}{+}\PY{l+m+mi}{2}\PY{o}{/}\PY{l+m+mi}{1000}\PY{p}{)}\PY{l+s+si}{:}\PY{l+s+s2}{.1f}\PY{l+s+si}{\PYZcb{}}\PY{l+s+s2}{\PYZpc{} of the time.}\PY{l+s+se}{\PYZbs{}n}\PY{l+s+s2}{\PYZdq{}}\PY{p}{)}
\end{Verbatim}
\end{tcolorbox}

    \begin{Verbatim}[commandchars=\\\{\}]

|   ANSWER 13
|
|   After the main sequence at 10 Gyr, stages 5 and 6 are the longest at 10 Myr;
|   each is smaller by a factor of 1000. As a fraction, this means that a star
is on
|   the main sequence 1/(1+2/1000) of the time, or percentage-wise
|   99.8\% of the time.

    \end{Verbatim}

    


    % Add a bibliography block to the postdoc
    
    
    
\end{document}
